\documentclass[conference,compsoc]{IEEEtran}
\usepackage[utf8]{inputenc}
\usepackage{xspace}
\usepackage{graphicx}
\usepackage{float}
\usepackage{subcaption}
\usepackage{pgfplots}
\usepackage{amsthm}
\usepackage{amsmath}
\usepackage{amssymb}
\usepackage{booktabs}
\usepackage[nocompress]{cite}
\usepackage[colorlinks=false, pdfborder={0 0 0}, hidelinks]{hyperref}
\pgfplotsset{compat=1.16}
\usepgfplotslibrary[colorbrewer]
\usetikzlibrary{pgfplots.colorbrewer}
\usetikzlibrary{calc}
\definecolor{color1}{HTML}{0868ac}
\definecolor{color2}{HTML}{43a2ca}
\definecolor{color3}{HTML}{7bccc4}
\definecolor{color4}{HTML}{a8ddb5}
\definecolor{color5}{HTML}{ccebc5}
\definecolor{myPurple}{HTML}{7B3294}
\definecolor{myOrange}{HTML}{E66100}
\definecolor{myGreen}{HTML}{1B9E77}
\definecolor{myRed}{HTML}{D7191C}

\newtheorem{proposition}{Proposition}

\newtheorem{definition}{Definition}

\begin{document}
\sloppy

\title{Biometric Authentication Mechanisms}

% \author{\IEEEauthorblockN{Dor Malka}
% \IEEEauthorblockA{Technion,\\
% Email: dm29@campus.technion.ac.il}
% \and
% \IEEEauthorblockN{Ittay Eyal}
% \IEEEauthorblockA{Technion,\\
% Email: ittay@technion.ac.il}}

% \renewcommand{\shortauthors}{Malka et al.}
\maketitle
\tableofcontents

\begin{abstract}
%%% 
%%% 
%%% 
%%% 
\emph{Biometric credentials} are prevalent in modern authentication: 
Sensors evaluate a measurement and output \emph{accept} if it is above a \emph{match threshold} or \emph{reject} otherwise. 
A large body of work improves the accuracy of biometric sensors, but the state-of-the-art uses a heuristic threshold and does not consider their integration into \emph{authentication mechanisms}. 
Recent works design authentication mechanisms to maximize the probability that the mechanism accepts a user and rejects an attacker, but treat biometric credentials as \emph{atomic} (perfect match or reject, like passwords). 

We model authentication mechanisms with both biometric and atomic credentials and prove that, in general, the optimal threshold differs from the common heuristic. 
Moreover, the optimal threshold depends on the integration with other credentials. 
More surprisingly, combining a biometric credential with an atomic one always improves the success probability over using the atomic credential alone. 
This holds even if the biometric credential is inaccurate.
Using real biometric datasets, we quantify the gain from optimizing the threshold across different mechanisms.

The implication is that biometric sensor vendors should change their interface so integrators can configure thresholds, resulting in better security of deployed authentication mechanisms.
\end{abstract}

\section{Introduction}
%%%
%%%
%%%
\newcommand{\FAR}{\ensuremath{\textit{FAR}}\xspace}
\newcommand{\FRR}{\ensuremath{\textit{FRR}}\xspace}
\newcommand{\EER}{\ensuremath{\textit{EER}}\xspace}
\newcommand{\ANDintegration}{\ensuremath{
P^{\text{AND}}_{\text{success}}(t)
=
\bigl(1-\FRR(t)\bigr)
\Bigl(
P^{1}_{\text{safe}}
+
P^{1}_{\text{leak}}\bigl(1-\FAR(t)\bigr)
\Bigr)
}\xspace}
\newcommand{\ORintegration}{\ensuremath{
    P^{\text{OR}}_{\text{success}}(t)
=
\bigl(1-\FAR(t)\bigr)
\Bigl(
P^{1}_{\text{safe}}
+
P^{1}_{\text{loss}}\bigl(1-\FRR(t)\bigr)
\Bigr)
}\xspace}
\newcommand{\Psuccess}{\ensuremath{
P_{\text{success}}(t)=
\begin{cases}
P^{1}_{\text{safe}}, 
& t \le a_1, \\[10pt]

(P^{1}_{\text{safe}} + P^{1}_{\text{leak}})
- P^{1}_{\text{leak}}\,\dfrac{a_2 - t}{a_2 - a_1}, 
& a_1 < t \le u_1, \\[14pt]

\dfrac{u_2 - t}{u_2 - u_1}
\left[
(P^{1}_{\text{safe}} + P^{1}_{\text{leak}})
- P^{1}_{\text{leak}}\,\dfrac{a_2 - t}{a_2 - a_1}
\right],
& u_1 < t \le a_2, \\[16pt]

\dfrac{u_2 - t}{u_2 - u_1}
(P^{1}_{\text{safe}} + P^{1}_{\text{leak}}),
& a_2 < t \le u_2, \\[12pt]

0, 
& t > u_2
\end{cases}
}\xspace}
\newcommand{\contprobabilities}{
\begin{align*}
P^{2}_{\text{safe}}(t) &= 1 - \FAR(t) - \FRR(t) + \FAR(t) \cdot \FRR(t),\\
P^{2}_{\text{loss}}(t) &= \FRR(t)\bigl(1 - \FAR(t)\bigr),\\
P^{2}_{\text{leak}}(t) &= \FAR(t)\bigl(1 - \FRR(t)\bigr),\\
P^{2}_{\text{theft}}(t) &= \FAR(t) \cdot \FRR(t).
\end{align*}
}
Biometric credentials are increasingly adopted in modern user devices. They rely on unique physical properties that make them difficult to forge~\cite{biometric_importance}. These properties have made them widely adopted. Broadly, biometric verification appears in~81\% of smartphones and~70\% of financial banking services, mainly through face or fingerprint verification~\cite{most_smartphones_biometric,financial_biometric} for access control. However, they were not always applicable~\cite{Anil1997} and therefore, credentials were typically \emph{atomic}, a credential was either presented correctly or not, evaluated according to a perfect match criterion.

Biometric credentials are evaluated by a sensor. The sensor obtains a~\emph{sample}, calculates a \emph{similarity score} between that sample and a~\emph{reference} stored in advance, and outputs~\emph{accept} if the score is above a set~\emph{match threshold} or \emph{reject} otherwise.

Previous work~$(\S\ref{sec:2})$ has conducted extensive analysis of biometric credentials in isolation, primarily focusing on reducing the \emph{false acceptance}~(\FAR) and \emph{false rejection}~(\FRR) \emph{rates}~\cite{Anil1997,prabhakar2003,Cappelli2006}. Sarkar et al.~\cite{sarkar2020} introduce the Equal Error Rate (\EER, where~$\FAR = \FRR$), a heuristic threshold of biometric sensors. But these works do not analyze the integration of biometric credentials into an authentication mechanism. Other works analyze authentication mechanisms as predicates of credential possession~\cite{maram2024,mouallem2024,Eyal2021,Abhishek2013}, but treat all credentials as atomic. 

In this work, we study how to design optimal mechanisms with biometric credentials. First, we present a refined definition of the authentication problem with biometric credentials~$(\S\ref{sec:3})$. The system comprises an authentication mechanism and two players: a~\emph{user} and an~\emph{attacker}. Each player interacts with the mechanism by presenting credentials.

 In practice, biometric sensors typically expose only this binary outcome~\cite{biometric_sensor_ibm,biometric_sensor_apple,biometric_sensor_ico,biometric_sensor_synaptics,biometric_sensor_authgear,biometric_sensor_faceheart} and do not allow the authentication mechanism to tune the match threshold. A \emph{genuine distribution}~$f_U(s)$ describes the score distribution of samples that belong to the legitimate user, whereas the \emph{impostor distribution}~$f_A(s)$ describes the score distribution of samples that belong to the attacker, as shown in Figure~\ref{fig:intro-fig}. These distributions establish the error metrics known as~\FAR and~\FRR. The~\FAR is the integral of the impostor distribution from the match threshold to the impostor distribution's maximum score~(red volume in Figure~\ref{fig:intro-fig}). Similarly, the~\FRR is the integral of the genuine distribution from its minimum score to the match threshold~(blue volume in Figure~\ref{fig:intro-fig}).

We examine distribution functions that approximate a biometric credential score distribution for both players.
Intuitively, genuine users are expected to produce higher matching scores than attackers. Therefore, the user score distribution is expected to be shifted toward higher matching scores compared to the attacker distribution. Otherwise, the attacker would be more likely to produce scores that lead to successful authentication, rendering the biometric credential ineffective. The \emph{state probabilities} of the biometric credential follow directly from the~\FAR and~\FRR metrics, which are functions of the match threshold. 

An optimal mechanism has the maximum probability of success, which depends on the number of credentials and their state probabilities. We start by optimizing a single biometric credential mechanism~$(\S\ref{sec:4})$. We derive closed-form expressions of the mechanism's success probability which, in a single credential mechanism, reduces to the probability that the credential is \emph{safe}, i.e., available only to the user. We identify the optimal threshold of the biometric credential and show that, in general, the optimal threshold differs from the commonly used~\EER threshold. This difference becomes more significant when the attacker is more likely to produce higher matching scores~$s$ and when the user and attacker distributions are asymmetric.

\begin{figure}[t]
\centering
\begin{tikzpicture}
\begin{axis}[
    scale only axis,
    width=0.8\columnwidth,
    height=0.50\columnwidth,
    xlabel={Matching Score},
    ylabel={Probability Density},
    xmin=0, xmax=1,
    ymin=0,
    xtick=\empty,
    ytick=\empty,
    grid=none,
    legend style={draw=black, fill=white},
    legend pos=north west,
]

% FRR shaded area: user rejected
\addplot[
    fill=blue,
    fill opacity=0.20,
    draw=none,
    forget plot
]
table [x=x, y=FRR] {figs/fig_2continuous/gaussian_far_frr_data.txt}
\closedcycle;

% FAR shaded area: attacker accepted
\addplot[
    fill=red,
    fill opacity=0.20,
    draw=none,
    forget plot
]
table [x=x, y=FAR] {figs/fig_2continuous/gaussian_far_frr_data.txt}
\closedcycle;

% User distribution
\addplot[blue, thick, mark=none]
table [x=x, y=user] {figs/fig_2continuous/gaussian_far_frr_data.txt};
\addlegendentry{User}

% Attacker distribution
\addplot[red, thick, mark=none]
table [x=x, y=attacker] {figs/fig_2continuous/gaussian_far_frr_data.txt};
\addlegendentry{Attacker}

% Threshold line
\addplot[black, dashed, thick, mark=none]
coordinates {(0.5,0) (0.5,4.2)};
\addlegendentry{Threshold}

\end{axis}
\end{tikzpicture}
\caption{User and Attacker score distributions with \FAR and \FRR areas appear in red and blue respectively.}
\label{fig:intro-fig}
\end{figure}

We extend our analysis to a more complex case that combines two credentials~$(\S\ref{sec:5})$. These may include two biometric credentials, as in airport or border-control systems~\cite{multimodal_biometrics}. For example, border control authorities may~(1)~Use cameras to capture a live traveler image and compare it with the data stored in the microchip inside the e-passport and~(2)~Use fingerprint verification to check for outstanding warrants. To evaluate a mechanism comprising two biometric credentials, we consider first and second order approximations of the biometric credential distributions and use a numerical optimization algorithm to find the optimal thresholds for each biometric credential. The optimal thresholds of the biometric credentials in such a mechanism are not, in general, the same as the optimal thresholds if used individually. This difference arises as we optimize each credential with respect to the entire mechanism.
To demonstrate the potential impact, we present a synthetic experiment of two sampled score distributions and observe a reduction of over~66\% in the failure rate, when comparing the~\EER and the optimal threshold. 

Another prevalent use case combines two credentials in one biometric and one atomic credential configuration, as in cloud service authentication. For example, Google~\cite{google_passkeys_2023} assumes (1)~Possession of a trusted device like a phone or computer (Atomic credential) and (2)~Authentication via biometric credentials~\cite{mahfouz2025}. We derive closed-form expressions for the mechanism's success probability and find that the optimal threshold depends on the properties of the atomic credential. Our symbolic analysis reveals that the mechanism's success can primarily depend on either the atomic credential or the biometric credential, depending on their properties.

Motivated by this observation, we identify the conditions under which the integration of a biometric credential with an atomic credential will always enhance the mechanism's success probability~$(\S\ref{sec:6})$. First, we show that when an atomic credential can be \emph{safe} or \emph{stolen} only it is not worthwhile to integrate, as the atomic credential mechanism dominates the mechanism success probability. Then, we introduce a necessary and sufficient condition for an \emph{asymptotically separable} biometric credential: for every score above or below a match threshold, the error rates converge to zero. Suppose the atomic credential has non-zero probabilities for the states of \emph{safe} and either \emph{loss} or \emph{leak}.
We prove that every \emph{asymptotically separable} biometric credential integrated with such an atomic credential, yields a mechanism whose success probability dominates that of the atomic credential mechanism. However, biometric credentials often exhibit a more common and weaker behavior~\cite{Slob2005,anil2004}. Thus, we introduce \emph{asymptotic dominance}: the genuine distribution dominates the impostor distribution by a factor~$c$ for every score above a specific threshold, or the impostor distribution dominates the genuine distribution by a factor~$c$ for every score below a specific threshold. We show that this is a sufficient condition.

To validate these theoretical conditions on real-world data, we quantitatively evaluate our framework using public biometric data~\cite{fvc2002,fvc2004} and a standard fingerprint verification algorithm~\cite{sourceafis}~$(\S\ref{sec:7})$. We use a dataset of fingerprints from multiple users and evaluate the resulting matching scores. To produce the user distribution, we match samples of the same user taken from the training set and the test set. We produce the attacker distribution by comparing the user's samples against those of other users. When calculating the matching scores produced when different users attempt to authenticate as one another, we observe that the algorithm produces low scores. Therefore, the attacker distribution spans a very narrow interval, resulting in a small overlapping area between the user and the attacker distributions and a very high success probability. Thus, no significant improvement occurs when comparing the mechanism's success under the~\EER threshold and the optimal threshold. This indicates why the~\EER threshold is sufficient when different users try to authenticate as one another.

However, since an attacker does not attempt to authenticate using their own biometric features, the relevant threat model involves attempts to physically impersonate a legitimate user. 
We conduct a spoofing attack evaluation using the LivDet database~\cite{LivDet2009,LivDet2011,LivDet2013,LivDet2015}, which includes both genuine fingerprint scans and corresponding spoofing attempts, produced by spoofing materials such as \textit{Gelatine}, \textit{Woodglue}, \textit{Latex} and~\textit{Ecoflex}. We match the spoofing samples of each user against the his geniune templates to obtain the imposter score distribution. These spoof samples produce matching scores that closely resemble those of the legitimate user. Integrating the biometric credential into a mechanism comprising an atomic credential, reduces the mechanism’s failure probability by~30-70\% compared to its failure probability under the \EER threshold. Furthermore, the results show a~50\% reduction in the mechanism’s failure probability compared to a mechanism based solely on an atomic credential, demonstrating the benefit of integrating the biometric credential.

In summary~$(\S\ref{sec:8})$, our contributions are both theoretical and practical: (1)~Formalization of the mechanism's success function composed of biometric and atomic credentials; (2)~Symbolic analysis of the optimal threshold for biometric and atomic authentication mechanisms. (3)~Characterization of the conditions under which biometric credential integration enhances the mechanism's success; (4)~Proof that any biometric credential, under a specified condition, always enhances the mechanism's success. (5)~Empirical analysis using real biometric data. (6)~Direct implication that biometric sensor vendors should allow configure thresholds for immediate security improvements.

\section{Related Work}\label{sec:2}
Authentication is a constant active and prominent research area, with prior work broadly categorized into two main aspects: credentials and mechanisms. Credentials are commonly manifested as passwords~\cite{password1979}, One-Time-Password (OTP)~\cite{otp2017}, physical devices~\cite{FIDO2021} or biometric features such as fingerprints, face recognition or iris~\cite{anil2004}. 

Extensive work has evaluated and improved the accuracy of biometric systems~\cite{Cappelli2006,Maio2002,Anil1997,LivDet2009,LivDet2011,LivDet2013,LivDet2015}. Prabhakar et al.~\cite{prabhakar2003} examine the tradeoff between two key error metrics of a biometric system, namely False Acceptance Rate~(\FAR) and False Rejection Rate~(\FRR). Sarkar et al.~\cite{sarkar2020} define the term Equal Error Rate~(\EER), the threshold where~\FAR and~\FRR are equal. They heuristically suggest it as the threshold for biometric credentials. However, these works do not define a success criterion for authentication mechanisms, nor do they explore how the presence of additional credentials can influence the optimal matching threshold of the biometric system. Hong et al.~\cite{hong1998} integrate fingerprint and face recognition, each one operates independently, using an \emph{AND mechanism}. They design a decision fusion protocol that makes a decision considering the~\FAR error metric only. However, they do not explore how one biometric credential affects the other for optimal usage, nor do they analyze the integration of biometric credentials along with other credential types within an authentication mechanism.

An authentication mechanism is a function that evaluates a set of credentials provided by a player~\cite{Abhishek2013} and output accept or reject~\cite{MattBishop2004,schneider_auth_people}. Recent work has explored authentication mechanism defined as predicate of credential possession. Maram et al.~\cite{maram2024} study interactive authentication mechanisms, where the user and attacker interact with the authentication mechanism. Mouallem et al.~\cite{mouallem2024} compare mechanisms' success and aim to identify mechanisms with maximal success when composed of a large number of homogeneous credentials. Eyal~\cite{Eyal2021} analyzes mechanisms modeled as Boolean functions of the credentials. However, these works are partially applicable to our setting. While we adopt their mechanism's success definition and the way they define the credential possession states, they model biometric credentials as atomic credentials neglecting the tuneable threshold.

%%%
%%%
%%%
\section{Model}\label{sec:3}
We formalize the authentication problem by specifying an execution and the system participants~$(\S\ref{sec:3.1})$, credential types and their characteritics~$(\S\ref{sec:3.2})$, and the authentication mechanism and its success~$(\S\ref{sec:3.3})$.

\subsection{System Participants}\label{sec:3.1}
The system comprises an authentication mechanism~$M$ and two players: a user~$U$ and an attacker~$A$. In an execution~$E$, a player~$P$ \emph{presents} a credential~$c_i$ to the mechanism by calling a random function \emph{present credentials} denoted by~$\textit{present}_P(i)$. We abstract away the concrete implementation of a credential and model only its presentation. The value of the~\emph{presentation function} determines whether a credential is available to a player. The entry~$\sigma_i^P=1$ means that~$P$ successfully presented credential~$c_i$, and~$\sigma_i^P=0$ means that it is not. At the beginning of an execution~$E$, the system is parameterized such by an \emph{allocation function}. A credential is successfully presented with probability randomized by the allocation function denoted by~$\textit{allocate}_P(i)$.

\begin{definition}(Availability Vector)
    An \emph{availability vector} specifies, for each credential, whether it is successfully presented by a player~$P$ or not. Denote by $\sigma^U$ \textit{and} $\sigma^A$ the availability vectors of the user and of the attacker, respectively.
    \begin{equation*}
    \sigma_i^P = \textit{present}_P(i)
    \end{equation*}
\end{definition}
The outcome of the presentation function determines whether~$\sigma_i^P=1$ or $\sigma_i^P=0$. A scenario~$\sigma$ is a tuple of availability vectors~$(\sigma^U, \sigma^A)$. 

\subsection{Credentials}\label{sec:3.2}
A credential may correspond to secret data, such as a password or one time password~(OTP); a physical object, such as a phone or security key; or a biometric property, such as a fingerprint or face. We model these credentials as either \emph{atomic} or \emph{biometric}. Each credential can be either:~$\textit{Safe}$: Only available to the user, $\textit{Lost}$: Not available to either player, $\textit{Leaked}$: Available to both players, or~$\textit{Stolen}$: Only available to the attacker.
\subsubsection*{3.2.1}\label{sec:3.2.1}
Passwords, OTPs and security keys are credential types evaluated according to a perfect match criterion, namely, atomic credentials.
\begin{definition} (Atomic Credential)~For an atomic credential the allocation function~$\textit{allocate}_P(i)$ randomly draws the state probabilities of the credential, namely~$P_{\text{safe}}$, $P_{\text{loss}}$, $P_{\text{leak}}$ and~$P_{\text{stolen}}$.
\end{definition}
When a player~$P$ presents an atomic credential~$c_i$, the presentation function ouputs~1 with probability~$P_{\text{safe}} + P_{\text{leak}}$ when allocated to the user and~$P_{\text{leak}} + P_{\text{stolen}}$ when allocated to the attacker.
\subsubsection*{3.2.2}\label{sec:3.2.2}
Alternatively, a credential may correspond to a biometric property, such as a fingerprint or face. The biometric sensor evaluates a measurement of the biometric property and produces a score~$s$. The user distribution~$f_U(s)$ denotes the score distribution of samples from the user, whereas the attacker distribution~$f_A(s)$ denotes the score distribution of samples from the attacker.
\begin{definition} (Biometric Credential)~For a biometric credential the presentation function~$\textit{present}_P(i)$ is defined by the user score distribution~$f_U(s)$, attacker score distribution~$f_A(s)$ and a match threshold~$t$. For a player~${P\in\{U,A\}}$, the allocation function~$\textit{allocate}_P(i)$ draws a random value for the score~$s$ from the distribution~$f_P(s)$. Thus
    \begin{equation*}
    s = \textit{allocate}_P(i)
    \end{equation*}
    \begin{equation*}
    \textit{present}_P(i) = \begin{cases}
        1 &  s \ge t \\
        0 & \text{otherwise}
    \end{cases}
    \end{equation*}
\end{definition}
We denote by~$\sigma_i$ the state of credential~$c_i$ in scenario~$\sigma$. 
% Each state of the biometric credential corresponds to the outcome of the presentation functions for a user and an attacker. Specifically, the pair $(\textit{present}_U(i), \textit{present}_A(i))$ determines whether the biometric credential is~\emph{Safe}, \emph{Loss}, \emph{Leak} or \emph{Stolen}. Therefore, the state probabilities are
% \begin{equation*}
% P_{\text{safe}} = \Pr[\mathsf{present}_U(i) = 1 \wedge \mathsf{present}_A(i) = 0].
% \end{equation*}
% \begin{equation*}
% P_{\text{loss}} = \Pr[\mathsf{present}_U(i) = 0 \wedge \mathsf{present}_A(i) = 0].
% \end{equation*}
% \begin{equation*}
% P_{\text{leak}} = \Pr[\mathsf{present}_U(i) = 1 \wedge \mathsf{present}_A(i) = 1].
% \end{equation*}
% \begin{equation*}
% P_{\text{stolen}} = \Pr[\mathsf{present}_U(i) = 0 \wedge \mathsf{present}_A(i) = 1].
% \end{equation*}
The user and attacker distributions establish two error metrics:(1)~\FAR, the false acceptance rate, is the probability that the sensor outputs \emph{accept} for an attacker sample. We calculate it by integrating the attacker distribution from the match threshold to the maximum score, thus
\begin{equation}
    \label{eqn:FAR}
\FAR(t) = \int_{t}^{\infty} f_A[s] ds
\end{equation}
(2)~\FRR, the false rejection rate, is the probability that the sensor outputs \emph{reject} for a user sample. We calculate it by integrating the user distribution from the minimum score to the match threshold, thus
\begin{equation}
    \label{eqn:FRR}
\FRR(t) = \int_{-\infty}^{t} f_U[s] ds
\end{equation}
We use~\FAR and~\FRR to find the state probabilities of the biometric credential. The user has access to the biometric credential when the output of the biometric sensor is \emph{accept}, which occurs with probability~$1-\FRR(t)$. The attacker has access to the biometric credential when the output of the biometric sensor is \emph{accept}, which occurs with probability~$\FAR(t)$. Thus, the probabilities of the four states are
\begin{equation*}
    P_{\text{loss}} = \FRR(t) \cdot (1 - \FAR(t)),
\end{equation*}
\begin{equation*}
P_{\text{leak}} = \FAR(t) \cdot (1 - \FRR(t)),
\end{equation*}
\begin{equation*}
P_{\text{stolen}} = \FRR(t) \cdot \FAR(t), \textit{ and }
\end{equation*}
\begin{equation}\label{eqn:safe_prob}
P_{{\text{safe}}} =1-\mathrm{P_{\text{loss}}}-\mathrm{P_{\text{leak}}}-\mathrm{P_{\text{stolen}}} = (1 - \FAR(t)) \cdot (1 - \FRR(t)).
\end{equation}

% For example, if~$\sigma_i = \text{leak}$ then the credential is available to both players, namely~$\sigma_i^U = \text{True}$ and~$\sigma_i^A = \text{True}$. If~$c_i$ is an atomic credential, we can deduce that~$x_{c_i^{\mathrm{atm}}}^{U}=1$ $\ \text{and}\ x_{c_i^{\mathrm{atm}}}^{A}=1$. In the case where~$c_i$ is a biometric credential, this implies~$s_{i}^U \ge t \ \text{and}\ s_{i}^A \ge t$. The probabilities that a user or an attacker satisfies an atomic credential can be characterized using password statistics from company databases~\cite{password_leak}. In contrast, the probabilities that a user or an attacker satisfies a biometric credential are characterized by the score distributions of the biometric credential.
\subsection{Mechanism}\label{sec:3.3}

An authentication mechanism~$M$ receives the~\emph{evaluation} of the presentation function~$\textit{present}_P(i)$ for each player~$P$ and outputs a \emph{decision}.
\begin{definition} (Authentication Mechanism) An authentication mechanism $M$ is a predicate that takes as input:
\begin{itemize}
     \item~$\sigma^P$ the availability vector of a player~$P$ in scenario~$\sigma$ (perhaps an empty vector, represented by~$\varnothing$).
\end{itemize}
and outputs either~1, meaning that the mechanism accepts the player or~0, meaning that the mechanism rejects the player.
\end{definition}

We are now ready to define the mechanism's success.
The mechanism is successful if it grants access to the user only. 
\begin{definition} (Mechanism Success) A mechanism~$M$ is \emph{successful} in a scenario~$\sigma$ if for a player~$P$ and its availability vector~$\sigma^P$, the authentication mechanism decision~$\mathrm{decide}_M(\sigma^P)$ equals~1 for the user only over an execution~$E$, that is
\begin{equation*}
    \mathrm{decide}_M(\sigma^U)=1 \wedge \mathrm{decide}_M(\sigma^A)=0.
\end{equation*}
Otherwise, the mechanism is \emph{unsuccessful}.
\end{definition} 
The mechanism's \emph{success probability} is thus
\begin{equation*}
P_{\text{Success}}(M) =
\Pr\big[\mathrm{decide}_M(\sigma^U)=1 \wedge \mathrm{decide}_M(\sigma^A)=0\big].
\end{equation*}
For example, in an \emph{AND mechanism}~\cite{Eyal2021}, the mechanism grants access to a player if and only if all credentials are available to the player; thus,~$\mathrm{decide}_M(\sigma^P)=1$ if ~$\sigma^P_i$ equals~1 for every credential~$c_i$. In an \emph{OR mechanism}, the mechanism grants access to a player if at least one credential is available to the player; thus,~$\mathrm{decide}_M(\sigma^P)=1$ if ~$\sigma^P_i$ equals~1 for at least one credential~$c_i$.
\section{Mechanisms}\label{sec:4}
We explore optimal mechanism success probability with a single biometric credential. We start with first~$(\S\ref{sec:4.1})$ and second~$(\S\ref{sec:4.2})$ order approximations.

\begin{figure*}[t]
\centering

\begin{subfigure}[t]{0.30\textwidth}
\centering
\begin{tikzpicture}
\begin{axis}[
    scale only axis,
    width=0.8\linewidth,
    height=0.5\linewidth,
    xlabel={Matching Score},
    ylabel={Probability Density},
    xmin=0, xmax=1,
    ymin=0, ymax=4.8,
    grid=major,
    xtick={0,0.15,0.40,0.60,0.85},
    xticklabels={$0$,$a_1$,$a_2$,$u_1$,$u_2$},
    ytick=\empty,
    legend to name=uniformthreecaseslegend,
    legend columns=2,
    legend style={draw=black, fill=white}
]
\addplot+[blue, thick, mark=none, const plot]
table [x=s, y=PU_case1] {figs/fig_uniform/uniform_three_cases.txt};
\addlegendentry{User}

\addplot+[red, thick, mark=none, const plot]
table [x=s, y=PA_case1] {figs/fig_uniform/uniform_three_cases.txt};
\addlegendentry{Attacker}
\end{axis}
\end{tikzpicture}
\caption{Separated uniform PDFs.}
\label{fig:uniform_case1_pdf}
\end{subfigure}
\begin{subfigure}[t]{0.30\textwidth}
\centering
\begin{tikzpicture}
\begin{axis}[
    scale only axis,
    width=0.8\linewidth,
    height=0.5\linewidth,
    xlabel={Matching Score},
    ylabel={},
    xmin=0, xmax=1,
    ymin=0, ymax=4.8,
    grid=major,
    xtick={0,0.20,0.40,0.65,0.85},
    xticklabels={$0$,$u_1$,$a_1$,$a_2$,$u_2$},
    ytick=\empty
]
\addplot+[blue, thick, mark=none, const plot]
table [x=s, y=PU_case2] {figs/fig_uniform/uniform_three_cases.txt};

\addplot+[red, thick, mark=none, const plot]
table [x=s, y=PA_case2] {figs/fig_uniform/uniform_three_cases.txt};
\end{axis}
\end{tikzpicture}
\caption{The attacker range lies within the user range.}
\label{fig:uniform_case2_pdf}
\end{subfigure}
\begin{subfigure}[t]{0.30\textwidth}
\centering
\begin{tikzpicture}
\begin{axis}[
    scale only axis,
    width=0.8\linewidth,
    height=0.5\linewidth,
    xlabel={Matching Score},
    ylabel={},
    xmin=0, xmax=1,
    ymin=0, ymax=4.8,
    grid=major,
    xtick={0,0.15,0.40,0.65,0.85},
    xticklabels={$0$,$a_1$,$u_1$,$a_2$,$u_2$},
    ytick=\empty
]
\addplot+[blue, thick, mark=none, const plot]
table [x=s, y=PU_case3] {figs/fig_uniform/uniform_three_cases.txt};

\addplot+[red, thick, mark=none, const plot]
table [x=s, y=PA_case3] {figs/fig_uniform/uniform_three_cases.txt};
\end{axis}
\end{tikzpicture}
\caption{Partial overlap of the distributions.}
\label{fig:uniform_case3_pdf}
\end{subfigure}

\ref*{uniformthreecaseslegend}

\caption{Quantitative illustration of three distinct cases of uniform attacker and user score distributions. In case (1), the distributions are separated. Case (2) illustrate the case where an attacker is fully contained in the user distribution, and (3) illustrate two overlapping configurations.}
\label{fig:uniform_three_cases}
\end{figure*}

\subsection{Uniform Probability Biometric Credentials}\label{sec:4.1}

We consider uniform distribution as a first-order approximation of the biometric credential distributions.
Denote by~$u_1 \textit{ and } u_2$ the user's PDF bounds and by~$a_1 \textit{ and } a_2$  the attacker's PDF bounds as shown in Figure~\ref{fig:uniform_three_cases}. The probability distribution functions for the user and the attacker, $f_U$ and $f_A$ are thus:
\begin{align*}
&f_U(s) = \frac{1}{u_2-u_1} \quad \text{for } u_1 \leq s \leq u_2 \text{ and}\\
&f_A(s) = \frac{1}{a_2-a_1} \quad \text{for } a_1 \leq s \leq a_2.
\end{align*}

The error metrics~\FAR and~\FRR are both functions of the match threshold~$t$ and given by Equations~\ref{eqn:FAR} and~\ref{eqn:FRR}. Thus:

\[
\FAR(t) = \int_{t}^{a_2} \frac{1}{a_2-a_1} ds =
\left\{
\begin{array}{ll}
1 & \text{if } t \le a_1 \\
\frac{a_2 - t}{a_2 - a_1} & \text{if } a_1 < t \le a_2 \\
0 & \text{if } t > a_2
\end{array}
\right.
\]
\[\textit{and}\]
\[
\FRR(t) = \int_{u_1}^{t} \frac{1}{u_2-u_1} ds =
\left\{
\begin{array}{ll}
0 & \text{if } t \le u_1 \\
\frac{t - u_1}{u_2 - u_1} & \text{if } u_1 < t \le u_2 \\
1 & \text{if } t > u_2
\end{array}
\right.
\]

For a biometric credentials to be useful, users are expected to produce higher matching scores than attackers. Otherwise, the attacker would be more likely to produce scores that lead to successful authentication, rendering the biometric credential ineffective. Thus, we assume that~${a_2 < u_2}$. There are three distinct cases: 

(1)~$a_2 < u_1$~(Figure~\ref{fig:uniform_case1_pdf}): In this case, the attacker and user distributions are \emph{separated}. For a mechanism with a single credential~$M^{\textit{single}}$, the mechanism's success probability is simply the probability that the credential is safe
\begin{equation}
    \label{eqn:safesuccess}
    P_{\text{success}}(M^{\text{single}}) = P_{{\text{safe}}}
\end{equation} 
Any threshold ~$a_2 \le t \le u_1$ yields zero~\FAR and~\FRR, and therefore achieves a mehcanism success probability of~$1$.

(2)~$u_1 < a_1 < a_2 < u_2$~(Figure~\ref{fig:uniform_case2_pdf}): The attacker distribution bounds lies within the range of the user distribution bounds. 
The success probability of the biometric credential is given by Equation~\ref{eqn:safe_prob}.
The~\FAR is one for any~$t \le a_1$ and zero for any~$t \ge a_2$.
Thus:
\begin{equation*}
P_{\text{success}}(t) = 
\begin{cases}
    0, & t \le a_1 \\
    (1 - \frac{a_2-t}{a_2-a_1})\cdot(1-\frac{t-u_1}{u_2-u_1}), & a_1 < t \le a_2\\
    1 - \frac{t-u_1}{u_2-u_1}, & a_2 < t \le u_2 \\
    0, & t > u_2
\end{cases}
\end{equation*}

(3)~$a_1 < u_1 < a_2 < u_2$~(Figure~\ref{fig:uniform_case3_pdf}): The user and attacker distributions partially overlap.

% \begin{figure}[t]
% \centering

% \begin{subfigure}{\columnwidth}
% \centering
% \begin{tikzpicture}
% \begin{axis}[
%     scale only axis,
%     width=0.65\columnwidth,
%     height=0.35\columnwidth,
%     title={Symmetric},
%     xlabel={Matching Score},
%     ylabel={Probability Density},
%     xmin=0, xmax=0.75,
%     ymin=0, ymax=3,
%     grid=major,
%     xtick={0,0.10,0.30,0.50,0.70},
%     xticklabels={$0$,$a_1$,$u_1$,$a_2$,$u_2$},
%     ytick={0,2.5},
%     yticklabels={$0$,$\frac{1}{u_2-u_1}$}
% ]
% \addplot+[blue, thick, mark=none, const plot]
% table [x=s, y=PU] {figs/fig_uniform/uniform_graphs_2.txt};

% \addplot+[red, thick, mark=none, const plot]
% table [x=s, y=PA_sym] {figs/fig_uniform/uniform_graphs_2.txt};
% \end{axis}
% \end{tikzpicture}
% \caption{Symmetric uniform PDFs.}
% \label{fig:uniform_symmetric_pdf}
% \end{subfigure}

% \begin{subfigure}{\columnwidth}
% \centering
% \begin{tikzpicture}
% \begin{axis}[
%     scale only axis,
%     width=0.65\columnwidth,
%     height=0.35\columnwidth,
%     title={Asymmetric},
%     xlabel={Matching Score},
%     ylabel={Probability Density},
%     xmin=0, xmax=0.75,
%     ymin=0, ymax=3.7,
%     grid=major,
%     xtick={0,0.05,0.30,0.42,0.70},
%     xticklabels={$0$,$a_1$,$u_1$,$a_2$,$u_2$},
%     ytick={0,2.5,3.3333},
%     yticklabels={$0$,$\frac{1}{u_2-u_1}$,$\frac{1}{a_2-a_1}$},
%     legend style={
%         at={(0.4,-0.32)},
%         anchor=north,
%         legend columns=2,
%         draw=black
%     }
% ]
% \addplot+[blue, thick, mark=none, const plot]
% table [x=s, y=PU] {figs/fig_uniform/uniform_graphs_2.txt};
% \addlegendentry{User}

% \addplot+[red, thick, mark=none, const plot]
% table [x=s, y=PA_asym] {figs/fig_uniform/uniform_graphs_2.txt};
% \addlegendentry{Attacker}
% \end{axis}
% \end{tikzpicture}
% \caption{Asymmetric uniform PDFs.}
% \label{fig:uniform_asymmetric_pdf}
% \end{subfigure}

% \caption{Symmetric and asymmetric uniform PDFs.}
% \label{fig:uniforms_distributions}
% \end{figure}

We maintain the same approach as the case~(2) to calculate the probability of success. The~\FAR is one for any~$t \le a_1$ and zero for any~$t \ge a_2$, and the~\FRR is zero for any~$t \le u_1$ and one for any~$t \ge u_2$. Thus, in this case, we obtain five regions of the success probability function:
\begin{equation*}
P_{\text{success}}(t) = 
\begin{cases}
    0, & t \le a_1 \\
    1 - \frac{a_2-t}{a_2-a_1}, & a_1 < t \le u_1 \\
    (1 - \frac{a_2-t}{a_2-a_1})\cdot(1-\frac{t-u_1}{u_2-u_1}), & u_1 < t \le a_2\\
    1 - \frac{t-u_1}{u_2-u_1}, & a_2 < t \le u_2 \\
    0, & t > u_2.
\end{cases}
\end{equation*}
We analyze the probability of success in each interval. First, we show that~$P_{\text{success}}(t)$ is continuous. Since~$P_{\text{success}}(t)$ is continuous on each region of its piecewise definition, it remains only to check the boundary points. At each boundary point, the left-hand and right-hand limits of~$P_{\text{success}}$ coincide with~$P_{\text{success}}$ itself; therefore,~$P_{\text{success}}$ is continuous on its entire domain. Then, we find the optimal threshold that maximizes the mechanism's success probability in each region. The second and third regions of case~(2) are equivalent to the third and fourth regions of case~(3), thus we analyze them together. In case~(3), the success probability function is linearly increasing in $a_1 \le t \le u_1$, and therefore its maximum is achieved at the right boundary,~$t = u_1$. In both cases, the success probability is linearly decreasing on $a_2 \le t \le u_2$, and therefore its maximum is achieved at the left boundary,~$t = a_2$. Since~$t = u_1$ and~$t = a_2$ are the boundary points of the region where~\FAR and~\FRR are strictly between zero and one, the maxima of the adjacent regions are already included in this region. Therefore, it is sufficient to analyze the region where~$u_1 \le t \le a_2$ in case~(3) or~$a_1 \le t \le a_2$ in case~(2),~\ref{fig:uniform_success_1_cred} illustrates this analysis.
Thus, we find the global maximum by differentiating the success probability function in this region and comparing the resulting critical point with the boundary values.

We differentiate the success function and set it to zero.
We find that at: $t = \frac{a_1+u_2}{2}$ we get $\frac{d}{dt} P_{\text{success}}(t) = 0$ with
$\frac{d}{dt} P_{\text{success}}(t) > 0 \text{, for every } t \in (u_1, \frac{a_1+u_2}{2})$
$\textit{and }$
$\frac{d}{dt} P_{\text{success}}(t) < 0 \text{, for every } t \in (\frac{a_1+u_2}{2}, a_2).$
Thus, the maximum point is at~$t_{\text{Max}} = \frac{a_1 + u_2}{2}$ if~$\frac{a_1+u_2}{2} \in [u_1, a_2]$ in case~(3) and~$\frac{a_1+u_2}{2} \in [a_1, a_2]$ in case~(2).

For both cases, the probability of success at the optimal point if the condition is satisfied is
\begin{equation*}
P_{\text{success}}(t_{\text{opt}}) = \frac{(a_1 - u_2)^2}{4 (a_1 - a_2) (u_1 - u_2)}.
\end{equation*}
For the overlap case, if the condition is not satisfied, the maximum point is at the boundary of the interval, which is either~${t = u_1}$ or~${t = a_2}$. Thus
\begin{equation*}
P_{\text{success}}(t_{\text{opt}}) = \max\left\{\frac{u_1 - a_1}{a_2 - a_1}, \frac{u_2 - a_2}{u_2 - u_1}\right\}.  
\end{equation*}
For the case where the attacker range lies within the user's range, if the condition is not satisfied the maximum point is at the lower boundary of the third interval, where~${t = a_2}$. Thus
\begin{equation*}
P_{\text{success}}(t_{\text{opt}}) = =\frac{u_2 - a_2}{u_2 - u_1}=.
\end{equation*}
The maximum point represents the optimal threshold of the mechanism. From the optimal success probability expression, we can see that when an attacker distribution spans a wider interval, i.e., when~$(a_2 - a_1)$ increases, the mechanism’s success probability may significantly decrease.
%Figure~\ref{fig:uniform_success_1_cred} illustrates the probability of success for a single biometric credential mechanism as a function of threshold~$t$ in the overlap case. The figure shows that in the second interval the success function is monotonically increasing, while in the fourth interval it is monotonically decreasing. Therefore, the optimal operating point is achieved in the interval where~\FAR and~\FRR are strictly between zero and one, which in this scenario is where~$u_1 \le t \le a_2$.
We calculate the~\EER point and show that it is not, in general, equal to the optimal point:
\begin{equation*}
t_{\EER} = \frac{a_1 \cdot u_1 - a_2 \cdot u_2}{a_1 - a_2 + u_1 - u_2}.
\end{equation*}
When the attacker's range lies withing the user's range, the optimal threshold lies in the boundry~${t = a_2}$ or inside the third interval~$t=\frac{a_1 + u_2}{2}$. These points are not equal to the~\EER and leads to different success probability. Similarly, in the overlap case, the optimal threshold lies in one of the boundries~${t = u_1}$,~${t = a_2}$ or inside the third interval where~$t=\frac{a_1 + u_2}{2}$. Thus, in both cases, the optimal threshold generally differs from the~\EER. However, when the user and attacker distributions are symmetric, i.e., when the intervals of each distribution are equal:~${a_2 - a_1 = u_2 - u_1}$, the~\EER expression coincides with the optimal threshold. Figure~\ref{fig:FAR_FRR_uniform_1_cred} illustrates this exception.
\begin{figure}[t]
\centering
\begin{tikzpicture}
\begin{axis}[
    scale only axis,
    width=0.8\columnwidth,
    height=0.42\columnwidth,
    xlabel={$t$},
    ylabel={$P_{\mathrm{success}}(t)$},
    grid=none,
    xmin=0, xmax=0.75,
    ymin=0, ymax=0.8,
    legend style={
        at={(0.98,0.98)},
        anchor=north east,
        fill=white,
        fill opacity=0.85,
        font=\scriptsize,
        text opacity=1,
        draw=black
    },
    xtick={0,0.05,0.30,0.375,0.42,0.70},
    xticklabels={$0$,$a_1$,$u_1$,$t_{\mathrm{opt}}$,$a_2$,$u_2$},
    ytick={0,0.7137},
    yticklabels={
        $0$
    }
]

\addplot[thick, myPurple, mark=none]
table [x=s, y=safe] {figs/fig_uniform/success_uniform_data.txt};

\addlegendentry{$P_{\mathrm{success}}(t)$}

\end{axis}
\end{tikzpicture}
\caption{Success probability as a function of the threshold.}
\label{fig:uniform_success_1_cred}
\end{figure}

\begin{figure}[t]
\centering
\begin{tikzpicture}
\begin{axis}[
    scale only axis,
    width=0.7\columnwidth,
    height=0.5\columnwidth,
    xlabel={$\FRR$},
    ylabel={$\FAR$},
    grid=major,
    xmin=-0.05, xmax=1.05,
    ymin=-0.05, ymax=1.05,
    xtick={0,0.2,0.4,0.6,0.8,1},
    ytick={0,0.5,1},
    legend style={
        at={(0.98,0.98)},
        anchor=north east,
        fill=white,
        fill opacity=0.85,
        font=\scriptsize,
        text opacity=1,
        draw=black
    }
]

\addplot[thick, myPurple, mark=none]
table [x=FRR, y=FAR] {figs/fig_uniform/success_uniform_data.txt};
\addlegendentry{Asymmetric Curve}

\addplot[thick, dashed, black, mark=none]
table [x=FRR_sym, y=FAR_sym] {figs/fig_uniform/success_uniform_data.txt};
\addlegendentry{Symmetric Curve}

\addplot[only marks, mark=*, mark size=1.3pt, blue]
table [x=FRR_opt, y=FAR_opt] {figs/fig_uniform/success_uniform_data.txt};
\addlegendentry{Optimal Point}

\addplot[only marks, mark=*, mark size=1.3pt, red]
table [x=FRR_eer, y=FAR_eer] {figs/fig_uniform/success_uniform_data.txt};
\addlegendentry{\EER Point}

\addplot[only marks, mark=*, mark size=1.3pt, blue]
table [x=FRR_opt_sym, y=FAR_opt_sym] {figs/fig_uniform/success_uniform_data.txt};


\addplot[only marks, mark=*, mark size=1.3pt, red]
table [x=FRR_eer_sym, y=FAR_eer_sym] {figs/fig_uniform/success_uniform_data.txt};

\end{axis}
\end{tikzpicture}
\caption{\FAR vs. \FRR operating curves for the asymmetric and symmetric distributions in~(3) case, with optimal and~\EER points.}
\label{fig:FAR_FRR_uniform_1_cred}
\end{figure}

\subsection{Quadratic Distributed Credential}\label{sec:4.2}
%%%
%%%
%%%
\begin{figure*}[t]
\centering

\begin{subfigure}[t]{0.30\textwidth}
\centering
\begin{tikzpicture}
\begin{axis}[
    scale only axis,
    width=0.8\linewidth,
    height=0.5\linewidth,
    xlabel={Matching Score},
    ylabel={Probability Density},
    xmin=0, xmax=1,
    ymin=0, ymax=7,
    grid=major,
    xtick={0,0.15,0.40,0.60,0.85},
    xticklabels={$0$,$\beta_1$,$\beta_2$,$\alpha_1$,$\alpha_2$},
    ytick=\empty,
    legend to name=quadraticthreecaseslegend,
    legend columns=2,
    legend style={draw=black, fill=white}
]
\addplot+[blue, thick, mark=none]
table [x=s, y=PU_case1] {figs/fig_parabola/quadratic_three_cases.txt};
\addlegendentry{User}

\addplot+[red, thick, mark=none]
table [x=s, y=PA_case1] {figs/fig_parabola/quadratic_three_cases.txt};
\addlegendentry{Attacker}
\end{axis}
\end{tikzpicture}
\caption{Separated quadratic PDFs.}
\label{fig:quadratic_case1_pdf}
\end{subfigure}
\begin{subfigure}[t]{0.30\textwidth}
\centering
\begin{tikzpicture}
\begin{axis}[
    scale only axis,
    width=0.8\linewidth,
    height=0.5\linewidth,
    xlabel={Matching Score},
    ylabel={},
    xmin=0, xmax=1,
    ymin=0, ymax=7,
    grid=major,
    xtick={0,0.20,0.40,0.65,0.85},
    xticklabels={$0$,$\alpha_1$,$\beta_1$,$\beta_2$,$\alpha_2$},
    ytick=\empty
]
\addplot+[blue, thick, mark=none]
table [x=s, y=PU_case2] {figs/fig_parabola/quadratic_three_cases.txt};

\addplot+[red, thick, mark=none]
table [x=s, y=PA_case2] {figs/fig_parabola/quadratic_three_cases.txt};
\end{axis}
\end{tikzpicture}
\caption{The attacker roots lies within the user roots.}
\label{fig:quadratic_case2_pdf}
\end{subfigure}
\begin{subfigure}[t]{0.30\textwidth}
\centering
\begin{tikzpicture}
\begin{axis}[
    scale only axis,
    width=0.8\linewidth,
    height=0.5\linewidth,
    xlabel={Matching Score},
    ylabel={},
    xmin=0, xmax=1,
    ymin=0, ymax=7,
    grid=major,
    xtick={0,0.15,0.40,0.65,0.85},
    xticklabels={$0$,$\beta_1$,$\alpha_1$,$\beta_2$,$\alpha_2$},
    ytick=\empty
]
\addplot+[blue, thick, mark=none]
table [x=s, y=PU_case3] {figs/fig_parabola/quadratic_three_cases.txt};

\addplot+[red, thick, mark=none]
table [x=s, y=PA_case3] {figs/fig_parabola/quadratic_three_cases.txt};
\end{axis}
\end{tikzpicture}
\caption{Partial overlap of the distributions.}
\label{fig:quadratic_case3_pdf}
\end{subfigure}

\ref*{quadraticthreecaseslegend}

\caption{Illustration of three distinct cases of quadratic attacker and user score distributions. In case (1), the distributions are separated. Case (2) illustrates the case where the attacker distribution is fully contained within the user distribution, and case (3) illustrates partial overlap between the two distributions.}
\label{fig:quadratic_three_cases}
\end{figure*}
Our goal is to model credentials whose distributions closely resemble those observed in real measurement data. To this end, we analyze a quadratic probability distribution function, which provides a second order approximation of the biometric credential that can also be seen in real fingerprint measurement data~$(\S\ref{sec:7})$. This choice is also consistent with the findings of Slobodan et al.~\cite{Slob2005} on Eigenfinger and Eigenpalm feature distributions. Accordingly, we define two probability density functions (PDFs). Let~$f_U$ denote the user distribution with roots~$\alpha_{1}, \alpha_{2}$ and let~$f_A$ denote the attacker distribution with roots~$\beta_{1}, \beta_{2}$. Donte by $a, b, c$ and $d, e, f$ the quadratic coefficients of the user and attacker distributions, respectively, hence:
\[
f_U(s) = 
\left\{
\begin{array}{ll}
    a \cdot s^2+b \cdot s+c & \alpha_1 \leq s \leq \alpha_2 \\
    0 & s > \alpha_2 \lor s < \alpha_1
    \end{array}
\right.
\]
\[
f_A(s) =
\left\{
\begin{array}{ll}
    d \cdot s^2+e \cdot s+f & \beta_1 \leq s \leq \beta_2 \\
    0 & s > \beta_2 \lor s < \beta_1
    \end{array}
\right.
\]

To reduce the degrees of freedom of the quadratic distributions, we express the coefficients directly in terms of their roots. The integral of each PDF over its roots interval is set to~1, which enforces the normalization condition. This condition unsures that the resulting error metrics~\FAR and~\FRR are bounded to~1. Under this condition, the coefficients become fully determined. Accordingly, we define the coefficients as
\begin{align*}
    &\alpha_1 \cdot \alpha_2 = \frac{c}{a} \rightarrow a = \frac{c}{\alpha_1 \cdot \alpha_2}\\
    &\alpha_1 + \alpha_2 = -\frac{b}{a} \rightarrow b = - \frac{c \cdot (\alpha_1 + \alpha_2)}{\alpha_1 \cdot \alpha_2} \\
    &\beta_1 \cdot \beta_2 = \frac{f}{d} \rightarrow d = \frac{f}{\beta_1 \cdot \beta_2}\\
    &\beta_1 + \beta_2 = -\frac{e}{d} \rightarrow e = - \frac{f \cdot (\beta_1 + \beta_2)}{\beta_1 \cdot \beta_2} \\
    &\int_{\alpha_1}^{\alpha_2} (a \cdot x^2+b \cdot x+c) dx = 1 \Rightarrow c = \frac{6 \cdot \alpha_1 \cdot \alpha_2}{(\alpha_1 - \alpha_2)^3} \\
    &\int_{\beta_1}^{\beta_2} (d \cdot x^2+e \cdot x+f) dx = 1 \Rightarrow f = \frac{6 \cdot \beta_1 \cdot \beta_2}{(\beta_1 - \beta_2)^3} \\
\end{align*}
Therefore, we can write the quadratic PDFs as
\begin{equation*}
f_U(s) =
\begin{cases}
     \frac{6s^2}{(\alpha_1 - \alpha_2)^3} - \frac{6s (\alpha_1 + \alpha_2)}{(\alpha_1 - \alpha_2)^3} + \frac{6 \alpha_1 \alpha_2}{(\alpha_1 - \alpha_2)^3} & \alpha_1 \leq s \leq \alpha_2 \\
    0, & \text{otherwise}
\end{cases}
\end{equation*}
\begin{equation*}
f_A(s) =
\begin{cases}
     \frac{6s^2}{(\beta_1 - \beta_2)^3} - \frac{6s  (\beta_1 + \beta_2)}{(\beta_1 - \beta_2)^3} + \frac{6  \beta_1 \cdot \beta_2}{(\beta_1 - \beta_2)^3} & \beta_1 \leq s \leq \beta_2 \\
    0, & \text{otherwise}
    \end{cases}
\end{equation*}
These distributions establish the error metrics,~\FAR and~\FRR, as a function of the match threshold~$t$. Denote
\begin{equation*}
A(t)=
\frac{
2(\beta_2^3-t^3)
-3(\beta_1+\beta_2)(\beta_2^2-t^2)
+6\beta_1\beta_2(\beta_2-t)
}{
(\beta_1-\beta_2)^3
},
\end{equation*}
\begin{equation*}
U(t)=
\frac{
2(t^3-\alpha_1^3)
-3(\alpha_1+\alpha_2)(t^2-\alpha_1^2)
+6\alpha_1\alpha_2(t-\alpha_1)
}{
(\alpha_1-\alpha_2)^3
}.
\end{equation*}

\begin{equation*}
\FAR(t)=
\begin{cases}
1, & t \le \beta_1,\\
A(t), & \beta_1 < t \le \beta_2,\\
0, & t > \beta_2,
\end{cases}
\end{equation*}
\[\textit{and}\]
\begin{equation*}
\FRR(t)=
\begin{cases}
0, & t \le \alpha_1,\\
U(t), & \alpha_1 < t \le \alpha_2,\\
1, & t > \alpha_2.
\end{cases}
\end{equation*}

Similarly to the analysis of the uniform distribution, the x axis represents the matching score $s$. Genuine users are expected to produce higher matching scores than attacker, thus we analyze three different cases:

(1)~$\beta_{2} < \alpha_{1}$~(Figure~\ref{fig:quadratic_case1_pdf}): The two distributions are separated. Any threshold $t$ chosen within the range~$\beta_{2} \le t \le \alpha_{1}$ yields zero~\FAR and~\FRR, and therefore achieves a mehcanism success probability of~$1$.

(2)~$\alpha_{1} < \beta_{1} < \beta_{2} < \alpha_{2}$~(Figure~\ref{fig:quadratic_case2_pdf}): The attacker distribution bounds lies within the range of the user distribution bounds. The mechanisms probability of success reduces to the safe probability of the biometric credential, which is given by Equation~\ref{eqn:safe_prob}
The~\FAR is one for any~$t \le \beta_1$ and zero for any~$t \ge \beta_2$. Thus
\begin{equation*}
\begin{aligned}
G(t)=
-&\frac{
(t-\alpha_2)^2(t-\beta_1)^2
}{
(\alpha_1-\alpha_2)^3(\beta_1-\beta_2)^3
} \\
&\cdot
(2t-3\alpha_1+\alpha_2)
(2t+\beta_1-3\beta_2).
\end{aligned}
\end{equation*}
\begin{multline*}
P_{\mathrm{success}}(t)=\\
\begin{cases}
0,
& t<\beta_1 \lor t>\alpha_2,\\
G(t),
& \beta_1 \le t \le \beta_2,\\
-\dfrac{(t-\alpha_2)^2(2t-3\alpha_1+\alpha_2)}
{(\alpha_1-\alpha_2)^3},
& \beta_2 \le t \le \alpha_2.
\end{cases}
\end{multline*}

(3)~$\beta_{1} \le \alpha_{1} \le \beta_{2} \le \alpha_{2}$~(Figure~\ref{fig:quadratic_case3_pdf}): The user and attacker distributions partially overlap. We mainatain same approach as the case~(2) to calculate the probability of success. The~\FAR is one for any~$t \le \beta_1$ and zero for any~$t \ge \beta_2$, and the~\FRR is zero for any~$t \le \alpha_1$ and one for any~$t \ge \alpha_2$. Thus, in this case, we obtain 4 regions of the success probability function
\begin{multline*}
P_{\mathrm{success}}(t)=\\
\begin{cases}
0,
& t<\beta_1 \lor t>\alpha_2,\\
\dfrac{(t-\beta_1)^2(2t+\beta_1-3\beta_2)}
{(\beta_1-\beta_2)^3},
& \beta_1 \le t \le \alpha_1,\\
G(t),
& \alpha_1 \le t \le \beta_2,\\
-\dfrac{(t-\alpha_2)^2(2t-3\alpha_1+\alpha_2)}
{(\alpha_1-\alpha_2)^3},
& \beta_2 \le t \le \alpha_2.
\end{cases}
\end{multline*}
The second and third intervals in case~(2) are identical to the third and fourth intervals in case~(3) and therefore, we analyze them together. 
We differentiate~$P_{\text{success}}(t)$ and compute its critical points in every interval. By comparing the stationary values obtained in all valid regions, we identify the global maximum, which corresponds to the optimal threshold~$t_{\text{opt}}$.
Before writing the piecewise derivative, we define the polynomial coefficients and the following function:
\begin{align*}
&q = 4, \\
&m = -(5\alpha_1 + \alpha_2 + \beta_1 + 5\beta_2),\\ 
&n = \alpha_1(3\alpha_2 + \beta_1 + 6\beta_2)
    - \alpha_2^2 + \alpha_2\beta_2 - \beta_1^2 + 3\beta_1\beta_2,\\
&\ell = \alpha_1\bigl(\beta_1^2 - 3\beta_2(\alpha_2 + \beta_1)\bigr)
       + \alpha_2^2 \beta_2,\\
&R(t) = -\frac{6 (t-\alpha_2)(t-\beta_1)\bigl(qt^3 + mt^2 + nt + \ell\bigr)}
     {(\alpha_1-\alpha_2)^3 (\beta_1-\beta_2)^3}.
\end{align*}
With these coefficients and function $R(t)$, the derivative of $P_{\text{success}}(t)$ is given by the following piecewise expression:
\begin{equation*}
\frac{dP_{\text{success}}}{dt} =
\begin{cases}
0, & t < \beta_1 \;\lor\; t > \alpha_2, \\

\dfrac{6 (t-\beta_1) (t-\beta_2)}{(\beta_1-\beta_2)^3},
& \beta_1 \le t \le \alpha_1, \\

R(t),
& \alpha_1 \le t \le \beta_2, \\

-\dfrac{6 (t-\alpha_1) (t-\alpha_2)}{(\alpha_1-\alpha_2)^3},
& \beta_2 \le t \le \alpha_2.
\end{cases}
\end{equation*}
In the second interval, the derivative has a single critical point at $t=\beta_{1}$. Since the derivative is strictly positive for all $t>\beta_{1}$, this point corresponds to a local minimum rather than a maximum. Similarly, in the fourth interval, the only critical point occurs at $t=\alpha_{2}$, and the derivative is strictly negative for all $t<\alpha_{2}$. Therefore, this interval also contains no maximum either. We therefore turn to the third interval. In this region, the derivative is the product of a cubic polynomial and two linear factors. The linear factors vanish at~$t=\beta_{1}$ and~$t=\alpha_{2}$, both of which lie outside the third interval. Consequently, the existence of critical points within this interval depends solely on the cubic term. Solving this cubic, reveals exactly one real root with the remaining two roots being complex. The closed form of the cubic solution cannot be simplified further more than a complex expression. Thus, we denote the solution as $t_{\text{cubic}}$.

\begin{figure*}[t]
\centering

\begin{subfigure}[t]{0.31\textwidth}
\centering
\begin{tikzpicture}
\begin{axis}[
    scale only axis,
    width=0.8\linewidth,
    height=0.45\linewidth,
    xlabel={$s$},
    ylabel={$P(s)$},
    xmin=1.5, xmax=4.5,
    ymin=-0.1, ymax=1.65,
    xtick={2,3,3.5,4},
    xticklabels={$\beta_1$,$\alpha_1$,$\beta_2$,$\alpha_2$},
    ytick=\empty,
    legend to name=distlegend,
    legend columns=2,
    legend style={draw=black, fill=white}
]
\addplot[blue, thick, mark=none]
table [x=s, y=PU] {figs/fig_parabola/parabola_case_a.txt};
\addlegendentry{User}

\addplot[red, thick, mark=none]
table [x=s, y=PA] {figs/fig_parabola/parabola_case_a.txt};
\addlegendentry{Attacker}
\end{axis}
\end{tikzpicture}
\caption{Attacker with low variance.}
\end{subfigure}
\hfill
\begin{subfigure}[t]{0.31\textwidth}
\centering
\begin{tikzpicture}
\begin{axis}[
    scale only axis,
    width=0.8\linewidth,
    height=0.45\linewidth,
    xlabel={$s$},
    ylabel={$P(s)$},
    xmin=1, xmax=4.5,
    ymin=-0.1, ymax=1.65,
    xtick={1.2,3,3.5,4},
    xticklabels={$\beta_1$,$\alpha_1$,$\beta_2$,$\alpha_2$},
    ytick=\empty
]
\addplot[blue, thick, mark=none]
table [x=s, y=PU] {figs/fig_parabola/parabola_case_b.txt};

\addplot[red, thick, mark=none]
table [x=s, y=PA] {figs/fig_parabola/parabola_case_b.txt};
\end{axis}
\end{tikzpicture}
\caption{Attacker with medium variance.}
\end{subfigure}
\hfill
\begin{subfigure}[t]{0.31\textwidth}
\centering
\begin{tikzpicture}
\begin{axis}[
    scale only axis,
    width=0.8\linewidth,
    height=0.45\linewidth,
    xlabel={$s$},
    ylabel={$P(s)$},
    xmin=0.5, xmax=4.5,
    ymin=-0.1, ymax=1.65,
    xtick={0.6,3,3.5,4},
    xticklabels={$\beta_1$,$\alpha_1$,$\beta_2$,$\alpha_2$},
    ytick=\empty
]
\addplot[blue, thick, mark=none]
table [x=s, y=PU] {figs/fig_parabola/parabola_case_c.txt};

\addplot[red, thick, mark=none]
table [x=s, y=PA] {figs/fig_parabola/parabola_case_c.txt};
\end{axis}
\end{tikzpicture}
\caption{Attacker with high variance.}
\end{subfigure}

\medskip

\pgfplotslegendfromname{distlegend}

% \Description{Three plots showing quadratic user and attacker probability density functions for three attacker configurations.}
\caption{Quadratic user and attacker distributions.}
\label{fig:different_attackers_pdf_1_quadratic_cred}
\end{figure*}

\begin{figure*}[t]
\centering

\begin{subfigure}[t]{0.31\textwidth}
\centering
\begin{tikzpicture}
\begin{axis}[
    scale only axis,
    width=0.8\linewidth,
    height=0.45\linewidth,
    xlabel={$t$},
    ylabel={$P_{\mathrm{success}}(t)$},
    xmin=1.5, xmax=4.5,
    ymin=0, ymax=1.05,
    xticklabels={$\beta_1$,$\alpha_1$,$\beta_2$,$\alpha_2$},
    xtick={2,3,3.5,4},
    ytick=\empty,
    legend to name=successlegend,
    legend columns=3,
    legend style={draw=black, fill=white}
]
\addplot[black, thick, mark=none]
table [x=s, y=Ps] {figs/fig_parabola/parabola_case_a.txt};
\addlegendentry{Success Curve}

\addplot[only marks, mark=*, mark size=1.3pt, blue]
table [x=T_opt, y=Ps_opt] {figs/fig_parabola/parabola_case_a_points.txt};
\addlegendentry{Optimal Point}

\addplot[only marks, mark=*, mark size=1.3pt, red]
table [x=T_eer, y=Ps_eer] {figs/fig_parabola/parabola_case_a_points.txt};
\addlegendentry{EER Point}
\end{axis}
\end{tikzpicture}
\caption{Attacker with low variance.}
\end{subfigure}
\hfill
\begin{subfigure}[t]{0.31\textwidth}
\centering
\begin{tikzpicture}
\begin{axis}[
    scale only axis,
    width=0.8\linewidth,
    height=0.45\linewidth,
    xlabel={$t$},
    ylabel={$P_{\mathrm{success}}(t)$},
    xticklabels={$\beta_1$,$\alpha_1$,$\beta_2$,$\alpha_2$},
    xtick={1.2,3,3.5,4},
    xmin=1, xmax=4.5,
    ymin=0, ymax=1.05,
    ytick=\empty
]
\addplot[black, thick, mark=none]
table [x=s, y=Ps] {figs/fig_parabola/parabola_case_b.txt};

\addplot[only marks, mark=*, mark size=1.3pt, blue]
table [x=T_opt, y=Ps_opt] {figs/fig_parabola/parabola_case_b_points.txt};

\addplot[only marks, mark=*, mark size=1.3pt, red]
table [x=T_eer, y=Ps_eer] {figs/fig_parabola/parabola_case_b_points.txt};
\end{axis}
\end{tikzpicture}
\caption{Attacker with medium variance.}
\end{subfigure}
\hfill
\begin{subfigure}[t]{0.31\textwidth}
\centering
\begin{tikzpicture}
\begin{axis}[
    scale only axis,
    width=0.8\linewidth,
    height=0.45\linewidth,
    xlabel={$t$},
    ylabel={$P_{\mathrm{success}}(t)$},
    xticklabels={$\beta_1$,$\alpha_1$,$\beta_2$,$\alpha_2$},
    xtick={0.6,3,3.5,4},
    xmin=0.5, xmax=4.5,
    ymin=0, ymax=1.05,
    ytick=\empty
]
\addplot[black, thick, mark=none]
table [x=s, y=Ps] {figs/fig_parabola/parabola_case_c.txt};

\addplot[only marks, mark=*, mark size=1.3pt, blue]
table [x=T_opt, y=Ps_opt] {figs/fig_parabola/parabola_case_c_points.txt};

\addplot[only marks, mark=*, mark size=1.3pt, red]
table [x=T_eer, y=Ps_eer] {figs/fig_parabola/parabola_case_c_points.txt};
\end{axis}
\end{tikzpicture}
\caption{Attacker with high variance.}
\end{subfigure}

\medskip

\pgfplotslegendfromname{successlegend}

% \Description{Three plots showing the success probability as a function of the threshold, with optimal and \EER points marked.}
\caption{Success probability for three quadratic attacker distributions.}
\label{fig:success_per_attacker_1_quadratic_cred}
\end{figure*}

\begin{figure*}[t]
\centering

\begin{subfigure}[t]{0.31\textwidth}
\centering
\begin{tikzpicture}
\begin{axis}[
    scale only axis,
    width=0.8\linewidth,
    height=0.45\linewidth,
    xlabel={$\FRR(t)$},
    ylabel={$\FAR(t)$},
    xmin=0, xmax=1,
    ymin=0, ymax=1,
    grid=none,
    ytick =\empty,
    xtick =\empty,
    legend to name=farfrrlegend,
    legend columns=3,
    legend style={draw=black, fill=white}
]
\addplot[thick, myPurple, mark=none]
table [x=FRR, y=FAR]
{figs/fig_parabola/parabola_case_a.txt};
\addlegendentry{Operating Curve}

\addplot[only marks, mark=*, mark size=1.3pt, blue]
table [x=FRR_opt, y=FAR_opt]
{figs/fig_parabola/parabola_case_a_points.txt};
\addlegendentry{Optimal Point}

\addplot[only marks, mark=*, mark size=1.3pt, red]
table [x=FRR_eer, y=FAR_eer]
{figs/fig_parabola/parabola_case_a_points.txt};
\addlegendentry{\EER Point}
\end{axis}
\end{tikzpicture}
\caption{Attacker with low variance.}
\end{subfigure}
\hfill
\begin{subfigure}[t]{0.31\textwidth}
\centering
\begin{tikzpicture}
\begin{axis}[
    scale only axis,
    width=0.8\linewidth,
    height=0.45\linewidth,
    xlabel={$\FRR(t)$},
    ylabel={$\FAR(t)$},
    xmin=0, xmax=1,
    ymin=0, ymax=1,
    grid=none,
    ytick =\empty,
    xtick =\empty
]
\addplot[thick, myPurple, mark=none]
table [x=FRR, y=FAR]
{figs/fig_parabola/parabola_case_b.txt};

\addplot[only marks, mark=*, mark size=1.3pt, blue]
table [x=FRR_opt, y=FAR_opt]
{figs/fig_parabola/parabola_case_b_points.txt};

\addplot[only marks, mark=*, mark size=1.3pt, red]
table [x=FRR_eer, y=FAR_eer]
{figs/fig_parabola/parabola_case_b_points.txt};
\end{axis}
\end{tikzpicture}
\caption{Attacker with medium variance.}
\end{subfigure}
\hfill
\begin{subfigure}[t]{0.31\textwidth}
\centering
\begin{tikzpicture}
\begin{axis}[
    scale only axis,
    width=0.8\linewidth,
    height=0.45\linewidth,
    xlabel={$\FRR(t)$},
    ylabel={$\FAR(t)$},
    xmin=0, xmax=1,
    ymin=0, ymax=1,
    grid=none,
    ytick =\empty,
    xtick =\empty
]
\addplot[thick, myPurple, mark=none]
table [x=FRR, y=FAR]
{figs/fig_parabola/parabola_case_c.txt};

\addplot[only marks, mark=*, mark size=1.3pt, blue]
table [x=FRR_opt, y=FAR_opt]
{figs/fig_parabola/parabola_case_c_points.txt};

\addplot[only marks, mark=*, mark size=1.3pt, red]
table [x=FRR_eer, y=FAR_eer]
{figs/fig_parabola/parabola_case_c_points.txt};
\end{axis}
\end{tikzpicture}
\caption{Attacker with high variance.}
\end{subfigure}

\medskip

\pgfplotslegendfromname{farfrrlegend}

% \Description{Three \FAR versus \FRR operating curves for different attacker distributions, each marking the optimal operating point and the equal error rate point.}
\caption{\FAR vs. \FRR operating curves for different attacker distributions.}
\label{fig:FAR_FRR_1_quadratic_cred}
\end{figure*}
By examining the sign of the derivative on both sides of the critical point $t_{\text{cubic}}$, we verify that it constitutes a global maximum of $P_{\text{safe}}(t)$. We denote this maximizing point by $t_{\text{opt}}$. Figure~\ref{fig:success_per_attacker_1_quadratic_cred} illustrates this analysis.

To demonstrate that this optimal threshold differs from the Equal Error Rate (\EER) threshold, we computed the~\EER point separately as follows:
\begin{equation*}
t_{\EER} = \frac{\alpha_1 \beta_1 - \alpha_2 \beta_2}{\alpha_1 - \alpha_2 + \beta_1 - \beta_2}
\end{equation*}
We can observe that the~\EER threshold is not equal to the~$t_{\text{opt}}$. To evaluate the difference between the~$t_{\text{opt}}$ and the~\EER, we illustrate different attacker distributions, each spanning a wider interval~$(\beta_1,\beta_2)$, which corresponds to increasing the variance, as shown in Figure~\ref{fig:different_attackers_pdf_1_quadratic_cred}.
Comparing each case with its corresponding success probability curve, we observe that the optimal threshold~$t_{\text{opt}}$ becomes more distinct from the~\EER threshold~$t_{\EER}$ as the variance of the attacker distribution increases and the more asymmetric the user and attacker distributions become. 

Moreover, we also observe that increasing the variance of an attacker distribution, shifts the optimal threshold~$t_{\text{opt}}$ towards~$\alpha_1$ reducing the~\FRR to zero, while the~\EER point remains fixed trying to balance between the~\FAR and~\FRR. This shift indicates that when increasing the variance, the optimal threshold prioritizes minimizing the false rejection rate, rather than balancing both error rates, as shown in Figure~\ref{fig:FAR_FRR_1_quadratic_cred}.

To quantify the effect of the variance on the deviation between the optimal threshold~$t_{\text{opt}}$ and~\EER threshold~$t_{\EER}$, we numerically computed this gap as a function of the attacker's variance. The simulation indicate that higher attacker variance enlarges the gap between the success probability at the optimal threshold and the success probability at the~\EER threshold as show in Figure~\ref{fig:gap_vs_variance_parabola_success}. This finding highlights the potential success advantage gained by operating at~$t_{\text{opt}}$ rather than at the~\EER point when an attacker distribution spans over wider intervals.

\begin{figure}[t]
\centering

% \begin{subfigure}[t]{\columnwidth}
% \centering
% \begin{tikzpicture}
% \begin{axis}[
%     scale only axis,
%     width=0.82\columnwidth,
%     height=0.48\columnwidth,
%     xlabel={Attacker Variance $\mathrm{Var}(s)$},
%     ylabel={$|t_{\mathrm{opt}}-t_{\EER}|$},
%     xtick=\empty,
%     ytick=\empty,
%     grid=none,
%     legend style={
%         at={(0.98,0.98)},
%         anchor=north east,
%         draw=black,
%         fill=white
%     }
% ]
% \addplot[thick, myPurple, mark=none]
% table [x=variance, y=gap] {figs/fig_parabola/gap_vs_variance_parabola.txt};
% \end{axis}
% \end{tikzpicture}
% \caption{Threshold gap.}
% \label{fig:gap_vs_variance_parabola_gap}
% \end{subfigure}

% \vspace{0.7em}

% \begin{subfigure}[t]{\columnwidth}
% \centering
\begin{tikzpicture}
\begin{axis}[
    scale only axis,
    width=0.82\columnwidth,
    height=0.48\columnwidth,
    xlabel={Attacker Variance $\mathrm{Var}(s)$},
    ylabel={$P_{\mathrm{success}}$},
    xtick=\empty,
    ytick=\empty,
    grid=none,
    legend style={
        at={(0.98,0.02)},
        anchor=south east,
        draw=black,
        fill=white
    }
]
\addplot[thick, blue, mark=none]
table [x=variance, y=Pmax] {figs/fig_parabola/gap_vs_variance_parabola.txt};
\addlegendentry{$P_{\max}$}

\addplot[thick, red, dashed, mark=none]
table [x=variance, y=PEER] {figs/fig_parabola/gap_vs_variance_parabola.txt};
\addlegendentry{$P_{\EER}$}
\end{axis}
\end{tikzpicture}
\caption{Success at the optimal and \EER thresholds.}
% \end{subfigure}

\caption{Comparison between the optimal threshold and the \EER threshold as a function of the attacker's variance.}
\label{fig:gap_vs_variance_parabola_success}
\end{figure}

\section{Two Credentials}\label{sec:5}
We explore mechanisms that integrate two independently and identically distributed credentials, derived from either the same or different types of credentials to model real life implications exist in cloud authentication for example. We numerically analyze two biometric credentials~$(\S\ref{sec:5.1})$ and symbolically analyze one biometric credential and one atomic credential~$(\S\ref{sec:5.2})$.

\subsection{Two Biometric Credentials}\label{sec:5.1}
%%%
%%%
To analyze the case of a mechanism $M$ comprises two biometric credentials, we conduct a numerical study over several probability distribution functions. For each distribution pair, we derive the optimal operating point for each possible mechanism composition by maximizing the respective success probabilities using a LBFGS optimization algorithm. We additionally compare the optimal operating point to the~\EER point, thereby characterizing the conditions under which the \emph{AND} or \emph{OR} mechanisms yields superior performance. 
We start by considering Gaussian distributions as a third-order approximation of a biometric credential's distribution. We defined the user and attacker distributions as~$f_U(s) = \mathcal{N}(\mu_u, \sigma_u^2)$ and~$f_A(s) = \mathcal{N}(\mu_a, \sigma_a^2)$ respectively. We randomly generate the Gaussian parameters to construct a user and an attacker PDF's closely resemble the empirical findings of Slobodan~\cite{Slob2005}. We focus on the overlapping case between the user and attacker distributions, as this represents the common case among real-world scenarios. For the first credential as~$\mu_u = 0.8$, $\sigma_u = 0.12$, $\mu_a = 0.3, \sigma_a = 1$ and for the second credential~$\mu_u = 1.0$, $\sigma_u = 0.2$, $\mu_a = 0.4, \sigma_a = 0.18$ as shown in Figure~\ref{fig:01_gaussians}.
\begin{figure}[t]
\centering

\begin{subfigure}{\columnwidth}
\centering
\begin{tikzpicture}
\begin{axis}[
    scale only axis,
    width=0.82\columnwidth,
    height=0.45\columnwidth,
    xlabel={$s$},
    ylabel={Density},
    grid=none,
    legend style={
        at={(0.98,0.98)},
        anchor=north east,
        draw=black,
        fill=white
    }
]

\addplot[blue, thick, mark=none]
table [x=s, y=PU1] {figs/fig_2continuous/gaussians_2cred_data.txt};
\addlegendentry{$f_U(s)$}

\addplot[red, thick, mark=none]
table [x=s, y=PA1] {figs/fig_2continuous/gaussians_2cred_data.txt};
\addlegendentry{$f_A(s)$}

\end{axis}
\end{tikzpicture}
\end{subfigure}

\begin{subfigure}{\columnwidth}
\centering
\begin{tikzpicture}
\begin{axis}[
    scale only axis,
    width=0.82\columnwidth,
    height=0.45\columnwidth,
    xlabel={$s$},
    ylabel={Density},
    grid=none,
]

\addplot[blue, thick, mark=none]
table [x=s, y=PU2] {figs/fig_2continuous/gaussians_2cred_data.txt};

\addplot[red, thick, mark=none]
table [x=s, y=PA2] {figs/fig_2continuous/gaussians_2cred_data.txt};

\end{axis}
\end{tikzpicture}
\end{subfigure}

% \Description{Two Gaussian density plots for two biometric credentials, each showing user and attacker score distributions.}
\caption{Gaussian user and attacker distributions.}
\label{fig:01_gaussians}
\end{figure}

For a mechanism $M$ comprises two independently and identically distributed credentials, derived from either the same or different PDF’s, we defined mechanism's success under two scenarios:

a. \emph{AND}: The user has availability to all credentials, while the attacker doesn't have an availability to at least one of the credentials. If both credentials are either in safe state, or one in safe state and the other in leak state the mechanism succeeds. If one of the credentials is either loss ot theft, a user~$U$ cannot satisfy the mechanism and thus the mechanism fails.
\begin{equation}
    \label{eqn:andsuccess}
P_{\text{Success}}(M^{\text{AND}}) = \mathrm{P^1_{safe}} \cdot \mathrm{P^2_{safe}} + \mathrm{P^1_{safe}} \cdot \mathrm{P^2_{leak}} + \mathrm{P^2_{safe}} \cdot \mathrm{P^1_{leak}}
\end{equation}

b. \emph{OR}: The user has availability to at least one credential, while the attacker has no availability to neither of them. If both credential are either in safe state, or one in safe state and the other in loss state the mechanism succeeds. If one of the credentials is either leak ot theft, an attacker can satisfy the mechanism and thus the mechanism fails.
\begin{equation}
    \label{eqn:orsuccess}
P_{\text{Success}}(M^{\text{OR}}) = \mathrm{P^1_{safe}} \cdot \mathrm{P^2_{safe}} + \mathrm{P^1_{safe}} \cdot \mathrm{P^2_{loss}} + \mathrm{P^2_{safe}} \cdot \mathrm{P^1_{loss}}
\end{equation}

We calculate the error metrics of each credential~$c_i$ which are given by Equations~\ref{eqn:FAR} and~\ref{eqn:FRR} as a function of threshold $t$ to calculate the state probabilities of each credential as seen in Section~\ref{sec:3.2}.
\[
\FAR(t) = \int_{-\infty}^{t} \frac{1}{\sigma_{a,i}\sqrt{2\pi}}
\exp\left(
-\frac{(s-\mu_{a,i})^2}{2\sigma_{a,i}^2}
\right)
ds, \textit{ and}
\]
\[
\FRR(t) = \int_{t}^{\infty} \frac{1}{\sigma_{u,i}\sqrt{2\pi}}
\exp\left(
-\frac{(s-\mu_{u,i})^2}{2\sigma_{u,i}^2}
\right)
ds.
\]

We analyze both~\emph{AND} and~\emph{OR} mechanisms success given by Equations~\ref{eqn:andsuccess} and~\ref{eqn:orsuccess} and find their optimal threshold~$t$, using LBFGS optimization algorithm. We also compute the~\EER point for each credential and find the corresponding success probability at this point.

The optimal operating point for both mechaninsms' types differ from the~\EER, yielding higher success probabilities than those obtained at the~\EER. Moreover, we find that the optimal points of both the \emph{AND} and \emph{OR} mechanisms are not necessarily the same as the optimal threshold if used individually. This finding suggest that integrators of a biometric credential should determine its threshold as function of the other credentials integrated into the mechanism. Figure~\ref{fig:02_far_vs_frr} illustrates the~\FAR vs. \FRR operating curves for both credentials, and marks the optimal points of both \emph{AND} and \emph{OR} mechanisms, as well as the~\EER and the standalone optimum of the biometric credential. The figure highlights the gap between different operating points. In the first credential, the threshold leads to high~\FRR and low~\FAR values, while in the second credential the threshold leads to low~\FRR and high~\FAR values. This demonstrates that the  \emph{AND} and \emph{OR} mechanisms balance the two error rates differently across credentials.

We also find that the integration of a biometric credential to a mechanism already comprises a biometric credential, improves both \emph{AND} and \emph{OR} mechanism success. Hence, integrating two biometric credentials into a mechanism improves its success over using them individually. Figure~\ref{fig:04_success_functions} illustrate the success probability of both \emph{AND} and \emph{OR} mechanisms. We fix a credential to its optimum and show how the success probability varies as a function of the thresholds of the other credential.

We numerically analyze first and second order approximations of the biometric credential's distribution and calculate the success probabilities of both \emph{AND} and \emph{OR} mechanisms at different thresholds. We quantify the improvement achieved by integrating two biometric credentials and the improvement of using the optimal threshold over the~\EER threshold. In conclusion, all numerical results are summarized in Table~\ref{tab:mechanism_success_gaussian}.

\begin{figure}[t]
\centering

\begin{subfigure}{\columnwidth}
\centering
\begin{tikzpicture}
\begin{axis}[
    scale only axis,
    width=0.82\columnwidth,
    height=0.45\columnwidth,
    xlabel={$\FAR$},
    ylabel={$\FRR$},
    xmin=0, xmax=0.5,
    ymin=0, ymax=0.5,
    grid=major,
    legend to name=farfrrtwocredlegend,
    legend columns=2,
    legend style={draw=black, fill=white}
]

\addplot[thick, myPurple, mark=none]
table [x=FAR1, y=FRR1] {figs/fig_2continuous/far_vs_frr_2cred_data.txt};
\addlegendentry{Operating Curve}

\addplot[only marks,  mark=*, mark size=1.3pt]
table [x=FAR1_and, y=FRR1_and] {figs/fig_2continuous/far_vs_frr_2cred_points.txt};
\addlegendentry{AND Optimum}

\addplot[only marks, myOrange, mark=x, mark size=1.8pt]
table [x=FAR1_or, y=FRR1_or] {figs/fig_2continuous/far_vs_frr_2cred_points.txt};
\addlegendentry{OR Optimum}

\addplot[only marks, myGreen, mark=square*, mark size=1.3pt]
table [x=FAR1_eer, y=FRR1_eer] {figs/fig_2continuous/far_vs_frr_2cred_points.txt};
\addlegendentry{\EER}

\addplot[only marks, myRed, mark=triangle*, mark size=1.5pt]
table [x=FAR1_safe, y=FRR1_safe] {figs/fig_2continuous/far_vs_frr_2cred_points.txt};
\addlegendentry{Safe-opt}

\end{axis}
\end{tikzpicture}
\end{subfigure}

\begin{subfigure}{\columnwidth}
\centering
\begin{tikzpicture}
\begin{axis}[
    scale only axis,
    width=0.82\columnwidth,
    height=0.45\columnwidth,
    xlabel={$\FAR$},
    ylabel={$\FRR$},
    xmin=0, xmax=0.5,
    ymin=0, ymax=0.5,
    grid=major
]

\addplot[thick, myPurple, mark=none]
table [x=FAR2, y=FRR2] {figs/fig_2continuous/far_vs_frr_2cred_data.txt};

\addplot[only marks, mark=*, mark size=1.3pt]
table [x=FAR2_and, y=FRR2_and] {figs/fig_2continuous/far_vs_frr_2cred_points.txt};

\addplot[only marks, myOrange, mark=x, mark size=1.8pt]
table [x=FAR2_or, y=FRR2_or] {figs/fig_2continuous/far_vs_frr_2cred_points.txt};

\addplot[only marks, myGreen, mark=square*, mark size=1.3pt]
table [x=FAR2_eer, y=FRR2_eer] {figs/fig_2continuous/far_vs_frr_2cred_points.txt};

\addplot[only marks, myRed, mark=triangle*, mark size=1.5pt]
table [x=FAR2_safe, y=FRR2_safe] {figs/fig_2continuous/far_vs_frr_2cred_points.txt};

\end{axis}
\end{tikzpicture}
\end{subfigure}

\medskip

\pgfplotslegendfromname{farfrrtwocredlegend}

% \Description{Two \FAR versus \FRR operating curves for two Gaussian credentials, with AND optimum, OR optimum, \EER, and standalone safe-opt points marked.}
\caption{\FAR vs. \FRR operating curves for two biometric credentials.}
\label{fig:02_far_vs_frr}
\end{figure}

\begin{figure}[t]
\centering

\pgfplotsset{
    successplot/.style={
        scale only axis,
        width=0.82\columnwidth,
        height=0.45\columnwidth,
        ylabel={$P_{\mathrm{success}}$},
        grid=major,
        each nth point=2,
        filter discard warning=false,
    }
}

\begin{subfigure}{\columnwidth}
\centering
\begin{tikzpicture}
\begin{axis}[
    successplot,
    xlabel={$T_1$},
    legend to name=success2credlegend,
    legend columns=2,
    legend style={draw=black, fill=white}
]

\addplot[thick, blue, mark=none]
table [x=T1, y=P_and_T1] {figs/fig_2continuous/success_functions_2cred_data.txt};
\addlegendentry{$P_{\mathrm{success}}^{\mathrm{AND}}$}

\addplot[thick, myOrange, mark=none]
table [x=T1, y=P_or_T1] {figs/fig_2continuous/success_functions_2cred_data.txt};
\addlegendentry{$P_{\mathrm{success}}^{\mathrm{OR}}$}

\addplot[only marks, mark=*, mark size=1.3pt, blue]
table [x=T1_and, y=P_and_opt] {figs/fig_2continuous/success_functions_2cred_points.txt};
\addlegendentry{AND Optimum}

\addplot[only marks, mark=x, mark size=1.8pt, myOrange]
table [x=T1_or, y=P_or_opt] {figs/fig_2continuous/success_functions_2cred_points.txt};
\addlegendentry{OR Optimum}

\addplot[only marks, mark=square*, mark size=1.3pt, myGreen]
table [x=T1_eer, y=P_and_T1_eer] {figs/fig_2continuous/success_functions_2cred_points.txt};
\addlegendentry{$P_{\mathrm{success}}^{\mathrm{AND}}(T_{\EER})$}

\addplot[only marks, mark=diamond*, mark size=1.5pt, myRed]
table [x=T1_eer, y=P_or_T1_eer] {figs/fig_2continuous/success_functions_2cred_points.txt};
\addlegendentry{$P_{\mathrm{success}}^{\mathrm{OR}}(T_{\EER})$}

\end{axis}
\end{tikzpicture}
\caption{Biometric credential 1 sweep.}
\end{subfigure}

\begin{subfigure}{\columnwidth}
\centering
\begin{tikzpicture}
\begin{axis}[
    successplot,
    xlabel={$T_2$}
]

\addplot[thick, blue, mark=none]
table [x=T2, y=P_and_T2] {figs/fig_2continuous/success_functions_2cred_data.txt};

\addplot[thick, myOrange, mark=none]
table [x=T2, y=P_or_T2] {figs/fig_2continuous/success_functions_2cred_data.txt};

\addplot[only marks, mark=*, mark size=1.3pt, blue]
table [x=T2_and, y=P_and_opt2] {figs/fig_2continuous/success_functions_2cred_points.txt};

\addplot[only marks, mark=x, mark size=1.8pt, myOrange]
table [x=T2_or, y=P_or_opt2] {figs/fig_2continuous/success_functions_2cred_points.txt};

\addplot[only marks, mark=square*, mark size=1.3pt, myGreen]
table [x=T2_eer, y=P_and_T2_eer] {figs/fig_2continuous/success_functions_2cred_points.txt};

\addplot[only marks, mark=diamond*, mark size=1.5pt, myRed]
table [x=T2_eer, y=P_or_T2_eer] {figs/fig_2continuous/success_functions_2cred_points.txt};

\end{axis}
\end{tikzpicture}
\caption{Biometric credential 2 sweep.}
\end{subfigure}

\medskip
\pgfplotslegendfromname{success2credlegend}

% \Description{Two success probability curves for threshold sweeps over two biometric credentials, showing AND and OR mechanisms with optimum and \EER thresholds.}
\caption{Success probability as a function of the threshold for two biometric credentials.}
\label{fig:04_success_functions}
\end{figure}

\begin{table*}[t]
\centering
\setlength{\tabcolsep}{4pt}
\begin{tabular}{lccccccccc}
\toprule
Distribution & $P_{c_1}(\mathrm{safe})$ & $P_{c_2}(\mathrm{safe})$ & $P_{\mathrm{AND}}(t_{\EER})$ & $P_{\mathrm{OR}}(t_{\EER})$ & $P_{\mathrm{AND}}(t_1^{\mathrm{safe}})$ & $P_{\mathrm{OR}}(t_2^{\mathrm{safe}})$ & $\mathbf{P_{\mathrm{AND}}(t_{\mathrm{opt}})}$ & $\mathbf{P_{\mathrm{OR}}(t_{\mathrm{opt}})}$ & Improvement \\
\midrule
Gaussian & 0.937 & 0.889 & 0.911 & 0.911 & 0.912 & 0.911 & \textbf{0.964} & \textbf{0.963} & 5.90\% / 5.66\% \\
Uniform & 1.000 & 0.916 & 0.956 & 0.956 & 0.917 & 0.999 & \textbf{1.000} & \textbf{1.000} & 4.63\% / 4.61\% \\
Parabolic & 0.959 & 0.887 & 0.921 & 0.921 & 0.920 & 0.923 & \textbf{0.987} & \textbf{0.985} & 7.25\% / 6.99\% \\
\bottomrule
\end{tabular}

\caption{Mechanism success under two biometric credentials with Gaussian, Uniform and Parabolic distributions at four reference operating points: $\mathrm{safe}$ success of a single biometric credential mechanism,~\EER-pair $(t_1^{\EER},t_2^{\EER})$, individual optimum pair $(t_1^{\mathrm{safe}},t_2^{\mathrm{safe}})$, and joint optima (AND at $(t_1^{*},t_2^{*})_{\mathrm{AND}}$ and OR at $(t_1^{*},t_2^{*})_{\mathrm{OR}}$)}. The improvement is measured between the optimal operating point and the~\EER point.
\label{tab:mechanism_success_gaussian}
\end{table*}

From~Table~\ref{tab:mechanism_success_gaussian}, we observe that the optimal operating point yields a higher success probability than the \EER point by between 5\% to 7\%. 

\subsection{One Atomic and One Biometric Credential}\label{sec:5.2}
%%%
%%%
%%%
We start by analyzing the case of an \emph{AND} mechanism comprises an atomic credential and a biometric credential. 
We define the atomic credential with a constant value of each probabilities' states as:~$P^1_{safe}$, $P^1_{loss}$, $P^1_{leak}$, $P^1_{theft}$ corrsponding respectively to each one of the states:~${\text{safe}}$,~${\text{loss}}$,~${\text{leak}}$ and~${\text{theft}}$. 
We define the biometric credential by its PDFs for the user and the attacker,~$f_U$ and~$f_A$ respectively. We consider uniform distribution as a first-order approximation of a biometric credential's distribution. Denote by~$u_1, u_2$ the user's PDF bounds and~$a_1, a_2$ the attacker's PDF bounds, as shown in Figure~\ref{fig:uniform_three_cases}. The state probabilities of the biometric credential are defined by the error metrics~\FAR and~\FRR as done in Section~\ref{sec:3.2}. Moreover, we consider the overlapping case between the user and attacker distributions, as this represents the common case among real-world scenarios. Thus, we assume that~$a_1 < u_1 < a_2 < u_2$.
We set the state probabilities of an atomic and a biometric credential into Equation~\ref{eqn:andsuccess}. Thus, the probability of success is a function of the threshold $t$ and divided into 5 intervals. Denote:
\[L(t) = \dfrac{u_2 - t}{u_2 - u_1}
\left[
P^{1}_{\mathrm{safe}} + P^{1}_{\mathrm{leak}}(1
- \dfrac{a_2 - t}{a_2 - a_1})
\right].\]
\[
P_{\mathrm{success}}(t)=
\begin{cases}
P^{1}_{\mathrm{safe}}, 
& t \le a_1 \\

P^{1}_{\mathrm{safe}} + P^{1}_{\mathrm{leak}}(1-\dfrac{a_2 - t}{a_2 - a_1}), 
& a_1 < t \le u_1 \\

L(t),
& u_1 < t \le a_2 \\

\dfrac{u_2 - t}{u_2 - u_1}\cdot(P^{1}_{\mathrm{safe}} + P^{1}_{\mathrm{leak}}),
& a_2 < t \le u_2 \\

0, 
& t > u_2 .
\end{cases}
\]

We analyze the success probability $P_{\text{success}}(t)$ by partitioning the threshold domain into five intervals.

The first interval, $t \le a_1$, corresponds to the case in which the threshold is lower than the minimum value of the attacker distribution. In this regime, the \FAR is equal to~1, while the \FRR is equal to~0. Consequently, the safe probability of the biometric credential is zero, and the success probability converges to $P^{1}_{\text{safe}}$. Since $P^{1}_{\text{safe}}$ is constant in this interval, no interior maximum exists.

The second interval, $a_1 < t \le u_1$, corresponds to the case in which the threshold lies between the minimum value of the attacker distribution and the minimum value of the user distribution. In this region, the \FAR decreases linearly as a function of~$t$, while the \FRR remains equal to~0. As a result, $P^2_{safe}$ and $P^2_{leak}$ of the biometric credential are $1-\FAR(t)$ and $\FAR(t)$, respectively. In this interval, $P_{\text{success}}(t)$ is a strictly increasing function, as its derivative is a constant positive value, and therefore the maximum is attained at the boundary point: $t_{\text{opt},2} = u_1$.


The third interval, $u_1 < t \le a_2$, corresponds to the case in which the threshold lies between the minimum value of the user distribution and the maximum value of the attacker distribution. In this region, both \FAR and \FRR vary linearly with~$t$, resulting in a success probability that is quadratic as a function of~$t$. Differentiating $P_{\text{success}}(t)$ yields a single critical point inside the interval,
\[
t_{\text{opt}}^\prime
=
\frac{
u_2 + a_2
- (a_2 - a_1)\,\dfrac{P^{1}_{\text{safe}} + P^{1}_{\text{leak}}}{P^{1}_{\text{leak}}}
}{2}
\]

Since this expression is valid only for the interval~$u_1 < t \le a_2$, the effective maximizer within this interval is obtained by clipping the optimum to the interval boundaries, as follows:
\[
t_{\text{opt},3}
=
\min\!\left(
a_2,\;
\max\!\left(u_1,\; t_{\text{opt}}^\prime\right)
\right).
\]
In particular, if~$t_{\text{opt},3} < u_1$, the maximum over this interval is  attained at~$t=u_1$, whereas if~$t_{\text{opt},3} > a_2$, the maximum is attained at~$t=a_2$. Evaluating the success probability at the user threshold~$t=u_1$ yields
\[
P_{\text{success}}(u_1)
=
P^{1}_{\text{safe}} + P^{1}_{\text{leak}}(1
- \frac{a_2 - u_1}{a_2 - a_1}).
\]
Since we are analyzing the overlapping case where~$u_1 > a_1$, it follows that
\[
P_{\text{success}}(u_1) > P^{1}_{\text{safe}},
\]
Demonstrating that the integration of the biometric credential strictly improves the success of the mechanism compared to using the atomic credential alone.

By examining the optimal point in this interval under \emph{AND} mechanism, we can deduce that as the variance of the attacker distribution increases, the optimal threshold $t_{\text{opt},3}$ shifts toward the lower boundry of the interval, $u_1$. Moreover, for an atomic credential with high safe values $P^{1}_{\text{safe}}$, and low leak values $P^{1}_{\text{leak}}$, the optimal threshold shifts toward the lower boundary as well. Thus, in such case, the success of the mechnaims becomes mainly dependent on the atomic credential. Conversely, an atomic credential with high leak values $P^{1}_{\text{leak}}$ and low safe values $P^{1}_{\text{safe}}$ causes the optimal threshold to shift toward the upper boundary of the interval, $a_2$. In this case, the success of the mechanism relies more heavily on the biometric credential.

The fourth interval,~$a_2 < t \le u_2$, corresponds to the case in which the threshold exceeds the maximum value of the attacker distribution. In this regime, the~\FAR is equal to~0, while the~\FRR increases linearly with~$t$. Consequently, $P_{\text{success}}(t)$ is a strictly decreasing function in this interval, and the maximum is attained at the boundary point
\[
t_{\text{opt},4} = a_2 .
\]

Finally, the fifth interval,~$t > u_2$, corresponds to the case in which the threshold exceeds the maximum value of the user and the attacker distributions, leaving the success probability equal to~0.

To determine the global optimal operating point, we compare all stationary points and boundary maxima obtained in the valid intervals. The optimal threshold is achieved in the third interval, where the attacker and user distributions overlap and the success function attains an interior maximum. Substituting $t_{\text{opt}}=t_{\text{opt},3}$ into $P_{\text{success}}(t)$ yields
\[
P_{\text{success}}(t_{\text{opt},3}) =
\frac{
\left(
- a_1 \big(P^{1}_{\text{safe}} + P^{1}_{\text{leak}}\big)
+ a_2 P^{1}_{\text{safe}}
+ P^{1}_{\text{leak}}\, u_2
\right)^2
}{
4 P^{1}_{\text{leak}} (a_2 - a_1)(u_2 - u_1)
}.
\]

We do similar analysis for the \emph{OR} mechanism. We find the probability of success for the \emph{OR} mechanism. Denote: 
\[
Q(t) = \dfrac{t - a_1}{a_2 - a_1}
\left[
P^{1}_{\mathrm{safe}} + P^{1}_{\mathrm{loss}}(1 - \dfrac{t - u_1}{u_2 - u_1})
\right].
\]
\[
P_{\text{success}}(t) =
\begin{cases}
0,
& t \le a_1 \\[8pt]

\dfrac{t - a_1}{a_2 - a_1}\,(P^{1}_{\mathrm{safe}} + P^{1}_{\mathrm{loss}}),
& a_1 < t \le u_1 \\[10pt]

Q(t),
& u_1 < t \le a_2 \\[12pt]

P^{1}_{\mathrm{safe}} + P^{1}_{\mathrm{loss}}(1 - \dfrac{t - u_1}{u_2 - u_1}),
& a_2 < t \le u_2 \\[10pt]

P^{1}_{\mathrm{safe}},
& t > u_2 .
\end{cases}
\]
We illustrate the mechanism's success probability for both \emph{OR} mechanism in figure~\ref{fig:OR_mechanism_success} and \emph{AND} mechanism in Figure~\ref{fig:AND_mechanism_success}. The figure marks the optimal operating point and the~\EER point on each curve demonstrating significant improvement from~\EER using optimal threshold. Moreover, the figure shows that when integrating a biometric credential and determining the threshold~$t$ to~\EER, the integration reduces mechanism's success compared to using the atomic credential alone. However, by optimizing the threshold, the integration of the biometric credential improves the mechanism's success over using the atomic credential alone.  
By differentiating every interval and identifying the critical points, we found that the optimal threshold that leads to the maximum probability of success lies in the third interval where~$u_1 < t \le a_2$:
\[t_{\text{opt}} = \frac{a_1\cdot{P^{1}_{\mathrm{loss}}} + u_2\cdot({P^{1}_{\mathrm{safe}}} + {P^{1}_{\mathrm{loss}}}) - P^{1}_{\mathrm{safe}}\cdot u_1}{2{P^{1}_{\mathrm{loss}}}}\]
And the maximum probability of success at this point is:
\[P_{\text{success}}(t_{\text{opt}}) = \frac{\big(a_1 {P^{1}_{\mathrm{loss}}} - u_2 ({P^{1}_{\mathrm{safe}}} + {P^{1}_{\mathrm{loss}}}) + {P^{1}_{\mathrm{safe}}} u_1\big)^2}
{4 {P^{1}_{\mathrm{loss}}} (a_2 - a_1)(u_2 - u_1)}\]

we calculate the~\EER point and show that the optimal operating point $t_{\text{opt}}$ differs from the equal error rate (\EER) point $t_{\EER}$.
\[t_{\EER} = \frac{a_1 u_1 - a_2 u_2}{a_1 - a_2 + u_1 - u_2} \neq t_{\text{opt}}.\]


\begin{figure}[t]
\centering

\pgfplotsset{
    mechsuccessplot/.style={
        scale only axis,
        width=0.75\columnwidth,
        height=0.5\columnwidth,
        xlabel={$t$},
        ylabel={$P_{\mathrm{success}}(t)$},
        ymin=-0.05,
        ymax=1.05,
        grid=major,
        each nth point=2,
        filter discard warning=false,
    }
}

% Read the points file
\pgfplotstableread{figs/fig_1continuous_1discrete/mechanism_1dis_1con_success_points.txt}{\successpoints}

% Extract OR values
\pgfplotstablegetelem{0}{P_opt_or}\of{\successpoints}
\edef\PoptOR{\pgfplotsretval}
\pgfplotstablegetelem{0}{P_or_eer}\of{\successpoints}
\edef\PeerOR{\pgfplotsretval}

% Extract AND values
\pgfplotstablegetelem{0}{P_opt_and}\of{\successpoints}
\edef\PoptAND{\pgfplotsretval}
\pgfplotstablegetelem{0}{P_and_eer}\of{\successpoints}
\edef\PeerAND{\pgfplotsretval}

\begin{subfigure}{\columnwidth}
\centering
\begin{tikzpicture}
\begin{axis}[
    mechsuccessplot,
    legend to name=biometricatomlegend,
    legend columns=4,
    legend style={draw=black, fill=white, font=\scriptsize},
    extra y ticks={\PeerOR,\PoptOR},
    extra y tick labels={$\PeerOR$,$\PoptOR$},
    extra y tick style={
        tick label style={font=\scriptsize},
        grid=major
    }
]

\addplot[thick, myOrange, mark=none]
table [x=T, y=P_or] {figs/fig_1continuous_1discrete/mechanism_1dis_1con_success_data.txt};
\addlegendentry{$P_{\mathrm{success}}^{\mathrm{OR}}(t)$}

\addlegendimage{thick, blue, mark=none}
\addlegendentry{$P_{\mathrm{success}}^{\mathrm{AND}}(t)$}

\addplot[only marks, mark=square*, mark size=1.3pt, myGreen]
table [x=T_eer, y=P_or_eer] {figs/fig_1continuous_1discrete/mechanism_1dis_1con_success_points.txt};
\addlegendentry{\EER}

\addplot[only marks, mark=*, mark size=1.5pt, black]
table [x=T_opt_or, y=P_opt_or] {figs/fig_1continuous_1discrete/mechanism_1dis_1con_success_points.txt};
\addlegendentry{Optimal point}

\end{axis}
\end{tikzpicture}
\caption{OR mechanism.}
\label{fig:OR_mechanism_success}
\end{subfigure}

\begin{subfigure}{\columnwidth}
\centering
\begin{tikzpicture}
\begin{axis}[
    mechsuccessplot,
    extra y ticks={\PeerAND,\PoptAND},
    extra y tick labels={$\PeerAND$,$\PoptAND$},
    extra y tick style={
        tick label style={font=\scriptsize},
        grid=major
    }
]

\addplot[thick, blue, mark=none]
table [x=T, y=P_and] {figs/fig_1continuous_1discrete/mechanism_1dis_1con_success_data.txt};

\addplot[only marks, mark=square*, mark size=1.3pt, myGreen]
table [x=T_eer, y=P_and_eer] {figs/fig_1continuous_1discrete/mechanism_1dis_1con_success_points.txt};

\addplot[only marks, mark=*, mark size=1.5pt, black]
table [x=T_opt_and, y=P_opt_and] {figs/fig_1continuous_1discrete/mechanism_1dis_1con_success_points.txt};

\end{axis}
\end{tikzpicture}
\caption{AND mechanism.}
\label{fig:AND_mechanism_success}
\end{subfigure}

\medskip
\pgfplotslegendfromname{biometricatomlegend}

\caption{Success probability as a function of the threshold for the OR and AND mechanisms composed of an atomic credential with state probabilities $P^{1}_{\mathrm{safe}} = 0.8$, $P^{1}_{\mathrm{loss}} = 0.05$, $P^{1}_{\mathrm{leak}} = 0.1$, $P^{1}_{\mathrm{theft}} = 0.05$, and a biometric credential.}
\label{fig:biometric_atomic_success}
\end{figure}

\section{Characterizing Success Improvement}\label{sec:6} 
\noindent
We characterize the conditions under which a mechanism that uses both a biometric credential and an atomic credential mechanism achieves greater success than a mechanism that uses only the atomic credential. We show that any mechanism using a single atomic credential with~${P^1_\text{loss} = P^1_\text{leak} = 0}$ and a biometric credential does not achieve greater success probability than a mechanism using only this atomic credential~$(\S\ref{sec:6.1})$. Next, we present a necessary and sufficient condition for such improvement~$(\S\ref{sec:6.2})$ and derive a more intuitive sufficient condition.
\subsection{No Improvement In Mechanism Success}\label{sec:6.1}
\noindent
We first consider an atomic credential that is always either safe or stolen, never lost or leaked. This case admits no benefit from producing a mechanism that uses an atomic credential in this setting and a biometric credential. The following proposition formalizes this result.
\begin{proposition}
    Let~$c_1$ be an atomic credential with~${P^1_\text{loss} = P^1_\text{leak} = 0}$, $c_2$ be a biometric credential with user distribution~$f_U(t)$ and attacker distribution~$f_A(t)$, and~$M$ be a mechanism using~$c_1$. Then, for a mechanism~$\bar{M}$ using~$c_1$,~$c_2$ in \emph{AND} or \emph{OR} compositions it holds that
    \begin{equation*}
        P_{\mathrm{success}}^{\bar{M}} \le P_{\mathrm{success}}^{M}.
    \end{equation*}
\end{proposition}
\begin{proof}
    The mechanism~$\bar{M}$ uses one atomic credential and one biometric credential under two possible compositions~(see~$\S\ref{sec:5}$). An \emph{AND mechanism} has success probability described by Equation~\ref{eqn:andsuccess}, and an \emph{OR mechanism} has success probability described by Equation~\ref{eqn:orsuccess}. For both equations,~$P^2_{\text{safe}}(t)$ denotes the safe probability of the biometric credential as a function of the operating threshold~$t$. Since the state probabilities of the atomic credential satisfy~$P^1_\text{leak} =  P^1_\text{loss} = 0$, the success probability of~$\bar{M}$ reduces in both cases to
    \begin{equation*}
        P_{\text{success}}(t) = P^1_\text{safe} \cdot P^2_{\text{safe}}(t).
    \end{equation*}
    Each state probability of the biometric credential is strictly less than one; in particular,~$P^2_{\text{safe}}(t) \le 1$ for all values of~$t$. Therefore,~$P_{\text{success}}(t) \le P^1_\text{safe}$ and thus, adding a biometric credential does not produce a new mechanism~$\bar{M}$ whose success probability is greater than that of a single atomic credential mechanism.
\end{proof}
\subsection{Necessary and Sufficient Condition for Integrating Biometric Credentials}\label{sec:6.2}
\noindent 
We turn to the general case of an atomic credential that can also be lost or leaked. Then, we identify a biometric credential that is \emph{asymptotically separable}. We show that this is a necessary and sufficient condition for a mechanism using an atomic credential and a biometric credential to achieve strictly greater success probability than a mechanism using only the atomic credential.
Intuitively, asymptotic separability means that there exists a threshold~$t$ such that the error rate is negligible with respect to the correctness rate, making the user and attacker distributions effectively separated.
\begin{definition}
A user distribution function~$f_U(s)$ and an attacker distribution function~$f_A(s)$ are \emph{asymptotically separable} if
\begin{equation*}
    \inf_{t}
    \min\!\left\{
    \frac{\FAR(t)}{1-\FRR(t)},\,
    \frac{\FRR(t)}{1-\FAR(t)}
    \right\}
    = 0.
\end{equation*}
\end{definition}
We now prove that asymptotic separability of a biometric credential is a sufficient condition.
\begin{proposition}\label{thm:suf_nec_cond_for_improve}
    Let~$c_1$ be an atomic credential with~$P^1_\text{safe} > 0$ and at least one of~$P^1_\text{leak},P^1_\text{loss} > 0$, $c_2$ be a biometric credential whose user distribution~$f_U(t)$ and attacker distribution~$f_A(t)$ are asymptotically separable and~$M$ be a mechanism using~$c_1$. Then, there exists a mechanism~$\bar{M}$ which is an \emph{AND} or \emph{OR} composition using~$c_1$ and~$c_2$, such that
    \begin{equation*}
        P_{\mathrm{success}}^{\bar{M}} > P_{\mathrm{success}}^{M}
    \end{equation*}
\end{proposition}

\begin{proof}
    The success probability of a mechanism comprising an atomic credential only is~$P^{1}_{\text{safe}}$~(Equation~\ref{eqn:safesuccess}). A biometric credential is characterized by four state probabilities, which are functions of~$t$: $P^{2}_{\text{safe}}(t)$, $P^{2}_{\text{loss}}(t)$, $P^{2}_{\text{leak}}(t)$, $P^{2}_{\text{theft}}(t)$. As previously presented in the model~$(\S\ref{sec:3.2})$, these probabilities are functions of~\FAR and~\FRR.
    \contprobabilities
    A mechanism~$\bar{M}$ using~$c_1$ and~$c_2$ has two possible compositions: an \emph{AND mechanism} with success probability described by Equation~\ref{eqn:andsuccess}, or an \emph{OR mechanism} with success probability described by Equation~\ref{eqn:orsuccess}. Thus, by placing the biometric credential in each configuration, we express the success probability of~$\bar{M}$ as a function of the operating threshold~$t$:
    \begin{align}
    &\ANDintegration \label{eqn:andineq}, \\
    &\ORintegration \label{eqn:orineq}.
    \end{align}
    We need to prove that there exists an operating threshold~$t$ such that
    $P^{\text{AND}}_{\text{success}}(t) > P^1_{\text{safe}}$ or $P^{\text{OR}}_{\text{success}}(t) > P^1_{\text{safe}}$.
    We plug Equation~\ref{eqn:andineq} into~$P^{\text{AND}}_{\text{success}}(t) > P^1_{\text{safe}}$ to achieve a sufficient condition for this inequality. Thus:
    \begin{align}
    &P^1_{\text{safe}} < (1-\FRR(t))\Big(P^1_{\text{safe}}+P^1_{\text{leak}}(1-\FAR(t))\Big) \iff \notag \\
    &\frac{P^{1}_{\text{safe}}}{P^{1}_{\text{safe}}+P^{1}_{\text{leak}}\cdot(1-\FAR(t))} < 1-\FRR(t)  \iff \notag \\
    &\FRR(t) < 1 - \frac{P^{1}_{\text{safe}}}{P^{1}_{\text{safe}}+P^{1}_{\text{leak}}\cdot(1-\FAR(t))} \iff \notag \\
    &\FRR(t) < \frac{P^{1}_{\text{leak}}\cdot(1-\FAR(t))}{P^{1}_{\text{safe}}+P^{1}_{\text{leak}}\cdot(1-\FAR(t))} \iff \notag \\
    &\frac{\FRR(t)}{1-\FAR(t)} < \frac{P^{1}_{\text{leak}}}{P^{1}_{\text{safe}}+P^{1}_{\text{leak}}\cdot(1-\FAR(t))}. \label{eqn:necessary term_1}
    \end{align}
    Division by~$(1-\FAR(t))$ in the fourth step is valid under the imposed constraint that~$P^{2}_\textit{safe}>0$, which implies~$\FAR<1$ and~$\FRR<1$.
    Since~$P^{1}_{\text{leak}}\cdot(1-\FAR(t)) \le P^{1}_{\text{leak}}$, a sufficient condition for the above inequality is
    \begin{equation}
        \label{eqn:FRRineq}
    \frac{\FRR(t)}{1-\FAR(t)} < \frac{P^{1}_{\text{leak}}}{P^{1}_{\text{safe}}+P^{1}_{\text{leak}}}
    \end{equation}

    Similarly, we plug Equation~\ref{eqn:orineq} into~$P^{\text{OR}}_{\text{success}}(t) > P^1_{\text{safe}}$, and deduce that this inequality holds when:
    \begin{equation}\label{eqn:necessary term_2}
        \frac{\FAR(t)}{1-\FRR(t)} < \frac{P^{1}_{\text{loss}}}{P^{1}_{\text{safe}}+P^{1}_{\text{loss}}\cdot(1-\FRR(t))}
    \end{equation}
    Division by~$(1-\FRR(t))$ is valid under the same imposed constraint that~$P^{2}_\textit{safe}>0$.
    Since $P^{1}_{\text{loss}}\cdot(1-\FRR(t)) \le P^{1}_{\text{loss}}$, a sufficient condition for the above inequality is
    \begin{equation}
        \label{eqn:FARineq}
        \frac{\FAR(t)}{1-\FRR(t)} < \frac{P^{1}_{\text{loss}}}{P^{1}_{\text{safe}}+P^{1}_{\text{loss}}}
    \end{equation}
    By the definition of asymptotic separability, there exists an operating threshold~$t$ such that 
    \begin{equation*}
    \inf_{t}
    \min\!\left\{
    \frac{\FAR(t)}{1-\FRR(t)},\,
    \frac{\FRR(t)}{1-\FAR(t)}
    \right\}
    = 0.
    \end{equation*}
    Equivalently, for all~$\varepsilon > 0$, there exists~$t$ such that either~$\frac{\FAR(t)}{1-\FRR(t)} < \varepsilon$ or~$\frac{\FRR(t)}{1-\FAR(t)} < \varepsilon$. Thus, we can choose~$\varepsilon$ that equals the right hand side of the sufficient condition where the \textsc{AND} mechanism is greater than~$M$~(Equation~\ref{eqn:FRRineq}) or the the right hand side where the~\textsc{OR} mechanism is greater than~$M$~(Equation~\ref{eqn:FARineq}). Hence, by definition of asymptotic separability we know that a~$t$ that satisfies the inequality exists. 
    More precisely, if \[\inf_{t}\frac{\FAR(t)}{1-\FRR(t)} = 0,\] 
    then the \textsc{OR} mechanism composition is preferable, whereas if \[\inf_{t}\frac{\FRR(t)}{1-\FAR(t)} = 0,\] then the \textsc{AND} mechanism composition is preferable. Therefore, a mechanism~$\bar{M}$ uses an asymptotically separable biometric credential and an atomic credential achieves a greater success probability than that of~$M$.
\end{proof}
We showed that \emph{asymptotic separability} is a sufficient condition. We now show it is a necessary condition for producing a new mechanism~$\bar{M}$ with greater success probability than that of~$M$.
\begin{proposition}\label{thm:nec_cond_for_improve}
    Let~$c_2$ be a biometric credential with user distribution~$f_U(t)$ and attacker distribution~$f_A(t)$. If for every atomic credential~$c_1$ with~$P^1_\text{safe} > 0$ and at least one of~$P^1_\text{leak},P^1_\text{loss} > 0$, there exists a threshold~$t$ and a mechanism~$\bar{M}$ formed as an \emph{AND} or \emph{OR} composition such that
    \begin{equation*}
    P^{\bar{M}}_{\text{success}}(t) > P^{M}_{\text{success}}.
    \end{equation*}
    Then the biometric credential~$c_2$ must be asymptotically separable.
\end{proposition} 
\begin{proof}
    Assume by contradiction that the biometric credential is not asymptotically separable. Then, there exists~$\varepsilon > 0$ such that for all~$t$ both~$\frac{\FAR(t)}{1-\FRR(t)} \ge \varepsilon$ and~$\frac{\FRR(t)}{1-\FAR(t)} \ge \varepsilon$. Since it is given that~$\bar{M}$ using both credentials~$c_1$ and~$c_2$ and achieves a strictly greater success probability than an atomic credential mechanism~$M$, then there exists a threshold~$t$ such that either:~$P^{\text{AND}}_{\text{success}}(t) > P^1_{\text{safe}} \textit{ or }$ $P^{\text{OR}}_{\text{success}}(t) > P^1_{\text{safe}}.$ Thus, for every atomic credential~$c_1$ necessarily either the inequality of \textsc{AND} mechanism holds~(Equation~\ref{eqn:necessary term_1}) or the inequality of \textsc{OR} mechanism holds~(Equation~\ref{eqn:necessary term_2}). In particular, these inequalities hold for an atomic credential with state probabilities~$P^{1}_{\text{safe}}, P^{1}_{\text{leak}}, P^{1}_{\text{loss}}$ such that
    \begin{equation*}
        \max\left\{\frac{P^{1}_{\text{leak}}}{P^{1}_{\text{safe}}}, \frac{P^{1}_{\text{loss}}}{P^{1}_{\text{safe}}}\right\} < \varepsilon.
    \end{equation*}
    Thus, for a mechanism~$\bar{M}$ comprising this atomic credential it holds that either
    \begin{equation*}
        \begin{aligned}
        &\frac{\FRR(t)}{1-\FAR(t)} < \frac{P^{1}_{\text{leak}}}{P^{1}_{\text{safe}}+P^{1}_{\text{leak}}\cdot(1-\FAR(t))} < \frac{P^{1}_{\text{leak}}}{P^{1}_{\text{safe}}} \textit{, or }\\
        &\frac{\FAR(t)}{1-\FRR(t)} < \frac{P^{1}_{\text{loss}}}{P^{1}_{\text{safe}}+P^{1}_{\text{loss}}\cdot(1-\FRR(t))} < \frac{P^{1}_{\text{loss}}}{P^{1}_{\text{safe}}}.
        \end{aligned}
    \end{equation*} 
    However, by the contradiction assumption that the biometric credential is not asymptotically separable it follows that for every threshold~$t$
    \begin{equation*}
        \begin{aligned}
            &\frac{\FRR(t)}{1-\FAR(t)} \ge \varepsilon > \frac{P^{1}_{\text{leak}}}{P^{1}_{\text{safe}}} \textit{, and } \\
            &\frac{\FAR(t)}{1-\FRR(t)} \ge \varepsilon > \frac{P^{1}_{\text{loss}}}{P^{1}_{\text{safe}}}.
        \end{aligned}
    \end{equation*} 
    This contradicts the fact that, for every atomic credential~$c_1$, the success probability of~$\bar{M}$ is strictly greater than that of~$M$, as presented in both Equation~\ref{eqn:necessary term_1} and Equation~\ref{eqn:necessary term_2}. Therefore, the biometric credential must be asymptotically separable.
\end{proof}
\noindent
We showed that~$\bar{M}$ achieves greater success probability than that of~$M$ if and only if the biometric credential is \emph{asymptotically separable} and the atomic credential has~$P^1_\text{safe} > 0$ and at least one of~$P^1_\text{leak},P^1_\text{loss} > 0$. We extend our proof to a mechanism that uses a biometric credential and~$k$ atomic credentials and show that the mechanism achieves greater success probability than that of a mechanism using only the atomic credentials. First, every authentication mechanism~$M$ can be represented as a sum of products, i.e., a disjunctive normal form~(DNF)~\cite{Mendelson1997}. Then, we define an induced credential of~$k$ atomic credentials as if it were itself a single atomic credential.

\begin{definition}
Let~$M$ be an authentication mechanism over~$k$ atomic credentials. The \emph{induced credential} of~$M$ is the single abstract credential whose availability to the user and attacker is determined by the outputs of~$M$ on their respective availability vectors~$\sigma^U,\sigma^A\in\{0,1\}^k$. The fault probabilities of the induced credential are thus
\begin{align*}
P^M_{\text{safe}}
&=
\Pr\left[M(\sigma^U)=1,\ M(\sigma^A)=0\right],\\
P^M_{\text{loss}}
&=
\Pr\left[M(\sigma^U)=0,\ M(\sigma^A)=0\right],\\
P^M_{\text{leak}}
&=
\Pr\left[M(\sigma^U)=1,\ M(\sigma^A)=1\right],\\
P^M_{\text{stolen}}
&=
\Pr\left[M(\sigma^U)=0,\ M(\sigma^A)=1\right].
\end{align*}
\end{definition}
\begin{proposition}\label{thm:suf_nec_cond_for_improve_k_atm}
    Let~$\{c_1,c_2,\ldots,c_k\}$ be a set of~$k$ independent atomic credentials,~$c_{k+1}$ be a biometric credential whose user distribution~$f_U(t)$ and attacker distribution~$f_A(t)$ are asymptotically separable and~$M$ be a mechanism using~$\{c_1,c_2,\ldots,c_k\}$ whose induced credential has~$P^M_\text{safe} > 0$ and at least one of~$P^M_\text{leak},P^M_\text{loss} > 0$. Then, there exists a mechanism~$\bar{M}$ using~$\{c_1,c_2,\ldots,c_k,c_{k+1}\}$ such that
    \begin{equation*}
        P_{\mathrm{success}}^{\bar{M}} > P_{\mathrm{success}}^{M}
    \end{equation*}
\end{proposition}
\begin{proof}
    Let~$M$ be a mechanism using~$\{c_1,c_2,\ldots,c_k\}$. The mechanism~$M$ induces a reduction of the~$k$ credentials into a single abstract credential. The induced credential of~$M$ is a single atomic credential with state probabilities~$P^M_\text{safe}>0$, and at least one of~$P^M_\text{leak}$, or~$P^M_\text{loss} > 0$. By Proposition~\ref{thm:suf_nec_cond_for_improve}, there exists a mechanism~$\bar{M}$ using the induced credential and the biometric credential that achieves greater success probability than that of the induced credential mechanism. Since the induced credential is equivalent to the original mechanism~$M$, it follows that there exists a mechanism~$\bar{M}$ using~$\{c_1,c_2,\ldots,c_k,c_{k+1}\}$ that achieves greater success probability than that of~$M$.
\end{proof}
\begin{proposition}
    Let~$c_{k+1}$ be a biometric credential with user distribution~$f_U(t)$ and attacker distribution~$f_A(t)$. Suppose that for every set of~$k$ independent atomic credentials~$\{c_1,\ldots,c_k\}$ and every mechanism~$M$ over these credentials whose induced credential has~$P^M_\text{safe} > 0$ and at least one of~$P^M_\text{leak},P^M_\text{loss} > 0$, there exists a mechanism~$\bar{M}$ using~$\{c_1,\ldots,c_k,c_{k+1}\}$ such that
    \begin{equation*}
    P^{\bar{M}}_{\text{success}}(t) > P^{M}_{\text{success}}.
    \end{equation*}
    Then the biometric credential~$c_{k+1}$ must be asymptotically separable.
\end{proposition} 
\begin{proof}
    The mechanism~$M$ induces a reduction of the~$k$ credentials into a single abstract credential. The induced credential of~$M$ is a single atomic credential with state probabilities~$P^M_\text{safe}>0$, and at least one of~$P^M_\text{leak},P^M_\text{loss} > 0$. Then,~$\bar{M}$ is a mechanism using the induced credential of~$M$ and the biometric credential. By Proposition~\ref{thm:nec_cond_for_improve}, the biometric credential must be asymptotically separable.
\end{proof}

Since biometric credentials often exhibit a more common behavior~\cite{Slob2005,anil2004}, we introduce the stronger \emph{asymptotic dominance} condition to capture this property. Intuitively, asymptotic dominance means that one distribution function dominates the other in the tails of the distributions: for very low score values~$s$ the attacker distribution dominates the user distribution, and vice versa for very high values of~$s$. We prove that a mechanism~$\bar{M}$ that uses an \emph{asymptotically dominated} biometric credential and an atomic credential achieves a greater success probability than that of a single atomic credential mechanism~$M$, and show it is a sufficient condition~(see Appendix~\ref{app:sufficient-condition}). 

\section{Measurements}\label{sec:7}
We now present experiments' results based on fingerprints taken from~$\textit{FVC2002}$ database~\cite{fvc2002} and \textit{FVC2004} database~\cite{fvc2004}. We use a standard fingerprint verification algorithm~\textit{SourceAFIS} created by Vazan~\cite{sourceafis} for scoring fingerprints out of their images based on minutiae triplets. We use fingerprints, as it is a widely common and used example of a biometric credential in authentication mechanisms. For a test set, we took five templates out of~$\textit{FVC2002}$ \textit{database1} of user 103 and compare them against the whole database of both~$\textit{FVC2002}$ and~$\textit{FVC2004}$. A geniune will be considered the templates of user 103 against its other templates while imposter considered as all the other combinations. For scoring methodology, we define the final score as the maximum similarity across all 5 user templates. For example, if the matching scores between a sample and the 5 probes are~$[0.1,0.3,0.6,0.9,0.4]$, then the resulting score is 0.9. 

We maintain same approach as Jain et al.~\cite{Anil1997}, Cappelli et al.~\cite{Cappelli2006} and Slobodan et al.~\cite{Slob2005} and build histograms of the geniune and imposter combinations as shown in Figure~\ref{fig:sourceafis_histogram}. Then, we turn the histogram into a smoothed probability distribution function. We find the~\FAR and~\FRR and calculate~\EER as shown in Figure~\ref{fig:sourceafis_far_frr}. 

The most important outcome from this approach was to achieve similar~\EER of around 2\% as achieved at industry and academical level sensors.~\cite{Maio2002}, to objectively make a comparison with state-of-the-art algorithms. We calculate the mechanisms successs probability as given in Equation~\ref{eqn:safesuccess} of a mechanism comprises a biometric credential and show that the optimal operating point is different from the~\EER point and achieves higher success probability as shown in Figure~\ref{fig:sourceafis_success}. Finally, we illustrate the probability success function of a mechanism composed of an atomic credential with~$P^{1}_\text{safe} = 0.9$ and~$P^{1}_\text{leak} = 0.04$, $P^{1}_\text{loss} = 0.04$ and $P^{1}_\text{stolen} = 0.02$ and integrate it with a biometric credential.

\begin{figure}[t]
\centering
\begin{tikzpicture}
\begin{axis}[
    ybar,
    bar width=3pt,
    width=0.9\columnwidth,
    height=0.65\columnwidth,
    xlabel={Normalized similarity score (\%)},
    ylabel={Count},
    grid=major,
    legend pos=north east,
    legend style={draw=black, fill=white, font=\scriptsize}
]

\addplot+[ybar, fill=blue, fill opacity=0.5]
table [x=bin, y=genuine] {figs/fig_different_users_sourceafsi/sourceafis_histogram_data.txt};
\addlegendentry{Genuine}

\addplot+[ybar, fill=red, fill opacity=0.5]
table [x=bin, y=impostor] {figs/fig_different_users_sourceafsi/sourceafis_histogram_data.txt};
\addlegendentry{Impostor}

\end{axis}
\end{tikzpicture}
\caption{Score histogram for genuine and impostor comparisons.}
\label{fig:sourceafis_histogram}
\end{figure}


\begin{figure}[t]
\centering
\begin{tikzpicture}
\begin{axis}[
    scale only axis,
    width=0.7\columnwidth,
    height=0.5\columnwidth,
    xlabel={Threshold (\%)},
    ylabel={Error Rate},
    ymin=0, ymax=1,
    grid=major,
    legend style={
    at={(0.5,-0.25)},
    anchor=north,
    legend columns=3,
    draw=black,
    fill=white,
    font=\scriptsize
    },  
    each nth point=2,
    filter discard warning=false
]

% FAR
\addplot[thick, red]
table [x=T, y=FAR] {figs/fig_different_users_sourceafsi/sourceafis_far_frr_data.txt};
\addlegendentry{FAR (smoothed-PDF integral)}

% FRR
\addplot[thick, blue]
table [x=T, y=FRR] {figs/fig_different_users_sourceafsi/sourceafis_far_frr_data.txt};
\addlegendentry{FRR (smoothed-PDF integral)}

% EER point
\addplot[only marks, mark=*, mark size=1.8pt, black]
table [x=T_eer, y=EER] {figs/fig_different_users_sourceafsi/sourceafis_far_frr_points.txt};
\addlegendentry{\EER}

\end{axis}
\end{tikzpicture}
\caption{\FAR and \FRR as a function of the decision threshold, with the Equal Error Rate (\EER) point.}
\label{fig:sourceafis_far_frr}
\end{figure}


\begin{figure}[t]
\centering

\pgfplotsset{
    sourceafissuccess/.style={
        scale only axis,
        width=0.8\columnwidth,
        height=0.50\columnwidth,
        xlabel={Threshold (\%)},
        ylabel={Success Probability},
        ymin=0,
        ymax=1,
        grid=major,
        legend style={
            at={(0.5,-0.25)},
            anchor=north,
            legend columns=3,
            draw=black,
            fill=white,
            font=\scriptsize
        },
        each nth point=2,
        filter discard warning=false
    }
}

\begin{subfigure}{\columnwidth}
\centering
\begin{tikzpicture}
\begin{axis}[
    sourceafissuccess
]

\addplot[thick, blue, mark=none]
table [x=T, y=P_success] {figs/fig_different_users_sourceafsi/sourceafis_p_success_data.txt};
\addlegendentry{$P_{\mathrm{success}}(t)$}

\addplot[only marks, mark=square*, mark size=1.4pt, myGreen]
table [x=T_eer, y=P_eer] {figs/fig_different_users_sourceafsi/sourceafis_p_success_points.txt};
\addlegendentry{$P_{\mathrm{success}}(t_{\EER})$}

\addplot[only marks, mark=*, mark size=1.6pt, myRed]
table [x=T_opt, y=P_opt] {figs/fig_different_users_sourceafsi/sourceafis_p_success_points.txt};
\addlegendentry{Optimal point}

\end{axis}
\end{tikzpicture}
\caption{Single biometric credential mechanism's success.}
\end{subfigure}

\begin{subfigure}{\columnwidth}
\centering
\begin{tikzpicture}
\begin{axis}[
    sourceafissuccess
]

\addplot[thick, blue, mark=none]
table [x=T, y=P_and] {figs/fig_different_users_sourceafsi/sourceafis_success_and_or_data.txt};
\addlegendentry{$P_{\mathrm{success}}^{\mathrm{AND}}(t)$}

\addplot[thick, myOrange, mark=none]
table [x=T, y=P_or] {figs/fig_different_users_sourceafsi/sourceafis_success_and_or_data.txt};
\addlegendentry{$P_{\mathrm{success}}^{\mathrm{OR}}(t)$}

\addplot[only marks, mark=*, mark size=1.5pt, blue]
table [x=T_and_opt, y=P_and_opt] {figs/fig_different_users_sourceafsi/sourceafis_success_and_or_points.txt};
\addlegendentry{AND optimum}

\addplot[only marks, mark=x, mark size=1.9pt, myOrange]
table [x=T_or_opt, y=P_or_opt] {figs/fig_different_users_sourceafsi/sourceafis_success_and_or_points.txt};
\addlegendentry{OR optimum}

\addplot[only marks, mark=square*, mark size=1.4pt, myGreen]
table [x=T_eer, y=P_and_eer] {figs/fig_different_users_sourceafsi/sourceafis_success_and_or_points.txt};
\addlegendentry{AND at \EER}

\addplot[only marks, mark=diamond*, mark size=1.5pt, myRed]
table [x=T_eer, y=P_or_eer] {figs/fig_different_users_sourceafsi/sourceafis_success_and_or_points.txt};
\addlegendentry{OR at \EER}

\end{axis}
\end{tikzpicture}
\caption{biometric credential integrated with an atomic credential into AND and OR mechanisms.}
\end{subfigure}

% \Description{Success probability curves as a function of threshold, including \EER and optimal operating points.}
\caption{
Success probability as a function of the threshold.
The AND and OR mechanisms achieve success probabilities of 
$0.93$ and $0.94$ at the \EER point, and 
$0.96$ and $0.95$ at their respective optimal operating points.
}
\label{fig:sourceafis_success}
\end{figure}

We show that the optimal operating point of the integrated mechanism is different from the \EER point and achieves higher success probability by around 3\% compared to the \EER point, as the \EER itself is only 1.8\%. In addition, we can observe that the optimal point in both mechanisms achieves higher success probability than both the atomic credetial mechanisms and biometric credential mechanism when tested individually as shown in Figure~\ref{fig:sourceafis_success}. Although there is an improvement in the success probability of the integrated mechanism, we can observe the improvement is not as significant as the one observed in the theoretical analysis. This can be explained by the fact that the \EER point of the biometric credential is already very low, indicating high seperability between a user and an attacker and thus, there is no much room for improvement. In that scenario, we compared several templates of user 103 against templates of different users. However, our main approach in this paper was to increase the mechanism's success in scenarios where an imposter will try to spoof the user templates and authenticate into the user account. Hence, we expect the imposter to use spoofing tools, producing samples resamble with score distribution function that much more similar to the genuine distribution function, which might lead to increased \EER point. 

\begin{figure}[t]
\centering
\begin{tikzpicture}
\begin{axis}[
    ybar,
    bar width=2pt,
    width=0.9\columnwidth,
    height=0.5\columnwidth,
    xlabel={Normalized similarity score (\%)},
    ylabel={Count},
    grid=major,
    legend pos=north east,
    legend style={draw=black, fill=white, font=\scriptsize},
]

\addplot[
    ybar,
    fill=blue,
    fill opacity=0.5,
    draw=blue
]
table [x=score, y=genuine] {figs/fig_spoofed_users/livdet_sourceafis_histogram_best_of_probes_data.txt};
\addlegendentry{Genuine}

\addplot[
    ybar,
    fill=red,
    fill opacity=0.5,
    draw=red
]
table [x=score, y=impostor] {figs/fig_spoofed_users/livdet_sourceafis_histogram_best_of_probes_data.txt};
\addlegendentry{Impostor}

\end{axis}
\end{tikzpicture}
% \Description{Histogram of normalized similarity scores for genuine and impostor comparisons using the best score over six probe templates.}
\caption{LivDet score histogram using the best score over probe templates.}
\label{fig:livdet_sourceafis_histogram}
\end{figure}


\begin{figure}[t]
\centering
\begin{tikzpicture}
\begin{axis}[
    scale only axis,
    width=0.8\columnwidth,
    height=0.5\columnwidth,
    xlabel={Threshold (\%)},
    ylabel={Error Rate},
    ymin=0,
    ymax=1,
    grid=major,
    legend pos=south east,
    legend style={legend columns=3,draw=black, fill=white, font=\scriptsize},
    each nth point=2,
    filter discard warning=false
]

\addplot[thick, red, mark=none]
table [x=T, y=FAR] {figs/fig_spoofed_users/livdet_sourceafis_far_frr_best_of_probes_data.txt};
\addlegendentry{FAR}

\addplot[thick, blue, mark=none]
table [x=T, y=FRR] {figs/fig_spoofed_users/livdet_sourceafis_far_frr_best_of_probes_data.txt};
\addlegendentry{FRR}

\addplot[only marks, mark=*, mark size=1.6pt, black]
table [x=T_eer, y=EER] {figs/fig_spoofed_users/livdet_sourceafis_far_frr_best_of_probes_points.txt};
\addlegendentry{\EER}

\end{axis}
\end{tikzpicture}
% \Description{\FAR and \FRR curves as a function of threshold for LivDet SourceAFIS scores using the best score over six probe templates, with the \EER point marked.}
\caption{\FAR and \FRR as a function of the threshold using the best score over probe templates. The Equal Error Rate (\EER) occurs at $\approx 17.8\%$.
}
\label{fig:livdet_sourceafis_far_frr}
\end{figure}


\begin{figure}[t]
\centering

\pgfplotsset{
    livdetsuccess/.style={
        scale only axis,
        width=0.8\columnwidth,
        height=0.50\columnwidth,
        xlabel={Threshold (\%)},
        ylabel={Success Probability},
        ymin=0,
        ymax=1,
        grid=major,
        legend pos=north east,
        legend style={
            legend columns=3,
            at={(0.5,-0.25)},
            anchor=north,
            draw=black, 
            fill=white, 
            font=\scriptsize},
        each nth point=2,
        filter discard warning=false
    }
}

\begin{subfigure}{\columnwidth}
\centering
\begin{tikzpicture}
\begin{axis}[livdetsuccess]

\addplot[thick, blue, mark=none]
table [x=T, y=P_success] {figs/fig_spoofed_users/livdet_sourceafis_p_success_best_of_probes_data.txt};
\addlegendentry{$P_{\mathrm{success}}(t)$}

\addplot[only marks, mark=square*, mark size=1.4pt, myGreen]
table [x=T_eer, y=P_eer] {figs/fig_spoofed_users/livdet_sourceafis_p_success_best_of_probes_points.txt};
\addlegendentry{$P_{\mathrm{success}}(t_{\EER})$}

\addplot[only marks, mark=*, mark size=1.6pt, myRed]
table [x=T_opt, y=P_opt] {figs/fig_spoofed_users/livdet_sourceafis_p_success_best_of_probes_points.txt};
\addlegendentry{Optimal point}

\end{axis}
\end{tikzpicture}
\caption{Single biometric credential.}
\end{subfigure}

\medskip

\begin{subfigure}{\columnwidth}
\centering
\begin{tikzpicture}
\begin{axis}[livdetsuccess]

\addplot[thick, blue, mark=none]
table [x=T, y=P_and] {figs/fig_spoofed_users/livdet_sourceafis_success_and_or_best_of_probes_data.txt};
\addlegendentry{$P_{\mathrm{success}}^{\mathrm{AND}}(t)$}

\addplot[thick, myOrange, mark=none]
table [x=T, y=P_or] {figs/fig_spoofed_users/livdet_sourceafis_success_and_or_best_of_probes_data.txt};
\addlegendentry{$P_{\mathrm{success}}^{\mathrm{OR}}(t)$}

\addplot[only marks, mark=*, mark size=1.5pt, blue]
table [x=T_and_opt, y=P_and_opt] {figs/fig_spoofed_users/livdet_sourceafis_success_and_or_best_of_probes_points.txt};
\addlegendentry{AND optimal point}

\addplot[only marks, mark=x, mark size=1.9pt, myOrange]
table [x=T_or_opt, y=P_or_opt] {figs/fig_spoofed_users/livdet_sourceafis_success_and_or_best_of_probes_points.txt};
\addlegendentry{OR optimal point}

\addplot[only marks, mark=square*, mark size=1.4pt, myGreen]
table [x=T_eer, y=P_and_eer] {figs/fig_spoofed_users/livdet_sourceafis_success_and_or_best_of_probes_points.txt};
\addlegendentry{AND at \EER}

\addplot[only marks, mark=diamond*, mark size=1.5pt, myRed]
table [x=T_eer, y=P_or_eer] {figs/fig_spoofed_users/livdet_sourceafis_success_and_or_best_of_probes_points.txt};
\addlegendentry{OR at \EER}

\end{axis}
\end{tikzpicture}
\caption{AND and OR mechanisms.}
\end{subfigure}

% \Description{Success probability as a function of threshold for the LivDet SourceAFIS scores using the best score over probe templates.}
\caption{
Success probability as a function of the decision threshold using the best score over probe templates.
The AND mechanism achieves a maximal success probability of $0.896$, while the OR mechanism achieves $0.826$.
At the \EER operating point, the success probabilities are $0.759$ and $0.684$ for the AND and OR mechanisms, respectively.
}
\label{fig:livdet_sourceafis_success}
\end{figure}
To support this claim, we use the LivDet2015 database \cite{LivDet2015,LivDet2013,LivDet2011,LivDet2009}, which contains a large-scale live and spoof fingerprint samples. 
We follow the same methodology as before and used the $\text{SourceAFIS}$ algorithm~\cite{sourceafis}. As probe templates, we select six Digital Persona scanner images of user 002\_0. Each probe was matched against 14 additional images of the same user to generate genuine scores. 
For each target image, we define the final score methodolodgy as the maximum similarity across 6 probe templates, similarly to the previous experiment. For example, if the matching scores between a target image and the six probes are [0.3, 0.5, 0.4, 0.8, 0.9, 0.7], then the resulting score is 0.9. 
To simulate a spoofing attack, we use fake fingerprint images of the same user (002\_0) produced by~\textit{Gelatine}, \textit{Woodglue}, \textit{Latex} and~\textit{Ecoflex} and match them against the probe templates to obtain the imposter score distribution. We present the resulting histogram in Figure~\ref{fig:livdet_sourceafis_histogram}. We then compute the~\EER point and plot the~\FAR-\FRR curve, as shown in Figure~\ref{fig:livdet_sourceafis_far_frr}. As expected, the \EER increased significantly to 17.8\%. Finally, we calculate the success probability for both a single biometric credential and for a mechanism comprises an atomic credential and a biometric credential, as shown in Figure~\ref{fig:livdet_sourceafis_success}. We can observe that the optimal operating point reducing the failure probability of an \textit{AND} mechanism by 60\% and \textit{OR} mechanism by 44\% compared to the~\EER point. Moreover, we reduced the failure probability by 50\% when integrating into \textit{AND} mechanism compared to an atomic credential mechanism with $P_{\text{safe}} = 0.8$, $P_{\text{loss}} = 0.04$, $P_{\text{leak}} = 0.15$, $P_{\text{theft}} = 0.01$ probabilities. We observe that fixing the match threshold~$t$ to the~\EER when we integrate the biometric credential into the mechanism, significantly reduces its success probability, compared to the success probability of an atomic credential mechanism only.

\section{Conclusion}\label{sec:8}
We introduce a formal model of authentication mechanisms comprising, for the first time, both biometric and atomic credentials. Our model characterizes each credential through four state probabilities, which allows us to analyze the success of different mechanism types. Using this model, we show that the optimal threshold of a biometric credential differs from the heuristic threshold commonly used today. We then symbolically calculate the optimal threshold and find that it depends on the other credentials integrated into the mechanism.

We prove that integrating an asymptotically separable biometric credential with an atomic one always improves the success probability, compared to a mechanism that comprises an atomic credential only.
We quantitatively demonstrate our theoretical analysis using public biometric datasets, simulating different scenarios, including spoofing attacks. Our results show that optimizing the threshold significantly reduces the failure probability when a spoofing attack occurs.

These findings suggest a new direction for designing authentication mechanisms. Instead of relying on a heuristic threshold, system designers should set the biometric match threshold according to the entire authentication mechanism. Developing new mechanism protocols based on this principle is a promising direction for future work.

Our findings also have important and impactful implications for biometric sensor vendors. Current biometric sensors often hide the threshold selection from mechanism designers and integrators, achieving a suboptimal mechanism success probability. We suggest that biometric sensor vendors should change their interface and let the integrators set the thresholds of the biometric credentials, based on the other credentials they integrate into the mechanism. Harnessing our approach may affect millions of users around the world by enhancing the security of their digital and physical assets.

{
\raggedright
\emergencystretch=3em
\Urlmuskip=0mu plus 2mu\relax
\def\UrlBreaks{%
  \do\a\do\b\do\c\do\d\do\e\do\f\do\g\do\h\do\i\do\j\do\k\do\l\do\m
  \do\n\do\o\do\p\do\q\do\r\do\s\do\t\do\u\do\v\do\w\do\x\do\y\do\z
  \do\A\do\B\do\C\do\D\do\E\do\F\do\G\do\H\do\I\do\J\do\K\do\L\do\M
  \do\N\do\O\do\P\do\Q\do\R\do\S\do\T\do\U\do\V\do\W\do\X\do\Y\do\Z
  \do\0\do\1\do\2\do\3\do\4\do\5\do\6\do\7\do\8\do\9
  \do\/\do\-\do\.\do\_\do\?\do\&\do\=\do\#}
\bibliographystyle{IEEEtran}
\bibliography{main}
}

\appendices

\section{A Sufficient Condition for Integrating Biometric Credentials}
\label{app:sufficient-condition}
\noindent
Since biometric credentials often exhibit a more common behavior~\cite{Slob2005,anil2004}, we introduce an \emph{asymptotic dominancy} condition which captures this property. We prove that adding an \emph{asymptotically dominated} biometric credential can produce a new mechanism~$\bar{M}$ with higher success probability than that of a single atomic credential mechanism~$M$, and show it is a sufficient condition.
\begin{definition}
\label{def:User and Attacker distribution functions}
A biometric credential with a user distribution function~$f_U(s)$ and an attacker distribution function~$f_A(s)$ are \emph{asymptotically dominated} if
\begin{equation*}
\begin{aligned}
&\forall c>0\;\exists t\; \text{s.t.}:\forall s \ge t: f_A(s)\le c\,f_U(s),\textit{ or} \\
&\forall c>0\;\exists t\; \text{s.t.}:\forall s \le t: f_U(s)\le c\,f_A(s).
\end{aligned}
\end{equation*}
\end{definition}

\begin{proposition}\label{thm:biometric-improves-security}
Let~$c_1$ be an atomic credential with~$P^1_\text{safe} > 0$ and at least one of~$P^1_\text{leak},P^1_\text{loss} > 0$, $c_2$ be a biometric credential whose user distribution~$f_U(t)$ and attacker distribution~$f_A(t)$ are asymptotically dominated and~$M$ be a mechanism using~$c_1$. Then, there exists a mechanism~$\bar{M}$ using~$c_1$,~$c_2$, such that
    \begin{equation*}
        P_{\mathrm{success}}^{\bar{M}} > P_{\mathrm{success}}^{M}
    \end{equation*}
\end{proposition}

\begin{proof}
The success probability of a mechanism comprising an atomic credential only is~$P^{1}_\textit{safe}$~(Equation~\ref{eqn:safesuccess}). A biometric credential is characterized by four state probabilities that are functions of~\FAR and ~\FRR
\contprobabilities
We previously observed that a sufficient condition for producing a new mechanism with greater success is described by Equation~\ref{eqn:FRRineq} and Equation~\ref{eqn:FARineq}. Thus, we need to show that there exists an operating threshold $t$ such that either
\begin{equation*}
    \begin{aligned}
    &\frac{\FRR(t)}{1-\FAR(t)} < \frac{P^{1}_{\text{leak}}}{P^{1}_{\text{safe}}+P^{1}_{\text{leak}}} \textit{, or } \\
    &\frac{\FAR(t)}{1-\FRR(t)} < \frac{P^{1}_{\text{loss}}}{P^{1}_{\text{safe}}+P^{1}_{\text{loss}}}
    \end{aligned}
\end{equation*}
We expand left hand side of Equation~\ref{eqn:FRRineq}. 
\begin{equation*}
\frac{\FRR(t)}{1-\FAR(t)} = \frac{\int_{-\infty}^{t} f_U(s) ds}{\int_{-\infty}^{t} f_A(s) ds}
\end{equation*}
By Definition~\ref{def:User and Attacker distribution functions}, since the biometric credential is asymptotically dominated, it follows that there exists an operating threshold~$t$ such that for every~$s \le t,$ $f_U(s) \le c\cdot\,f_A(s)$. Therefore:
\begin{equation*}
\frac{\int_{-\infty}^{t} f_U(s) ds}{\int_{-\infty}^{t} f_A(s) ds} \le \frac{\int_{-\infty}^{t} c \cdot f_A(s) ds}{\int_{-\infty}^{t} f_A(s) ds} = c
\end{equation*}
This holds~$\forall c > 0,$; thus, we can choose~$c$ to equal the right hand side of~Equation~\ref{eqn:FRRineq} and satisfy the sufficient condition.
Similarly, we expand left hand side of Equation~\ref{eqn:FARineq}. 
\begin{equation*}
\frac{\FAR(t)}{1-\FRR(t)} = \frac{\int_{t}^{\infty} f_A(s) ds}{\int_{t}^{\infty} f_U(s) ds}
\end{equation*}
By Definition~\ref{def:User and Attacker distribution functions}, since the biometric credential is asymptotically dominated, it follows that there exists an operating threshold~$t$ such that for every~$s \ge t,$ $f_A(s) \le c\cdot\,f_U(s)$. Therefore,
\begin{equation*}
\frac{\int_{t}^{\infty} f_A(s) ds}{\int_{t}^{\infty} f_U(s) ds} \le \frac{\int_{t}^{\infty} c \cdot f_U(s) ds}{\int_{t}^{\infty} f_U(s) ds} = c, \forall c > 0
\end{equation*}
We have showed that~$\forall c>0$ either~$\frac{\FAR(t)}{1-\FRR(t)} < c$ or $\frac{\FRR(t)}{1-\FAR(t)} < c$. Therefore, integrating an asymptotically dominated biometric credential with an atomic credential produces~$\bar{M}$ with greater success probability.
\end{proof}

This proof is non trivial, as the integration of an atomic credential mechanism with another atomic credential, does not necessarily produce a new mechanism with higher success probability. Moreover, we have shown that integrating a biometric credential into an atomic credential mechanism without adjusting the match threshold~$t$ may reduce the mechanism's success probability~$(\S\ref{sec:5.2})$. Thus, Proposition~\ref{thm:suf_nec_cond_for_improve} and Proposition~\ref{thm:biometric-improves-security} emphasize the importance of adjusting the threshold when adding a biometric credential into an authentication mechanism.
The \emph{asymptotically dominated} condition reflects practical scenarios, as for very low score values~$s$ the attacker distribution dominates the user distribution, and vice versa for very high values of~$s$. We validate this claim with real public biometric data in Section~\ref{sec:7}.

\end{document}
\endinput
